\begin{textsample}{POS Dim 3 – 2017\_c6 – Score 1.00 – 2017\_c6/enade\_2017\_qui\_5a144fb539.txt}  \label{ex:f3_pos_199}
CAPÍTULO 6 ANÁLISE TÉCNICA DA PROVA Este capítulo tem por objetivo apresentar o desempenho dos estudantes concluintes de Química ( Licenciatura/Bacharelado ) no Enade/2017.
Para isso, foram calculadas as estatísticas básicas da prova em seu todo ( seção 6.1.1 ), bem como as estatísticas relacionadas aos Componentes de Formação Geral ( seção 6.1.2 ) e de Conhecimento Específico da Área ( seção 6.1.3 ).
Dadas as suas características, foram analisadas, em separado, as questões objetivas ( seção 6.2 ) e as questões discursivas ( seção 6.3 ).
Para as questões objetivas, foram disponibilizados os índices de facilidade e de discriminação ponto-bisserial, também em separado, para os Componentes de Formação Geral ( seções 6.2.1 e 6.2.2 ) e de Conhecimento Específico ( seções 6.2.3 e 6.2.4 ).
De cada componente, uma das questões foi escolhida para exemplificar a análise gráfica, relacionando as alternativas escolhidas pelos estudantes ( inclusive o gabarito ) com o número de acertos no componente.
O Anexo I apresenta a íntegra da análise gráfica para todas as questões objetivas.
Para cada uma das questões discursivas, os conteúdos dos tipos mais comuns de respostas dos estudantes são apresentados e comparados com o padrão de respostas esperado ( ver Anexo VIII com o padrão de respostas ).
Tomando como base as duas questões discursivas do Componente de Formação Geral, a seção 6.3.1.6 apresenta comentários sobre a correção das respostas com respeito à Língua Portuguesa.
As tabelas apresentam as seguintes estatísticas das notas20 : média do desempenho na prova, erro padrão da média, desvio padrão, nota mínima, mediana e nota máxima para cada um de seus componentes.
Tais estatísticas contemplam o total de estudantes concluintes da Área de Química ( Licenciatura/Bacharelado ) inscritos e presentes à prova do Enade/2017, tendo em vista agregações, ou por Grandes Regiões e o país como um todo, ou por Categoria Administrativa, Organização Acadêmica da IES e Modalidade de Ensino.
Em relação aos gráficos de distribuição de notas, o intervalo considerado foi de 10 unidades, aberto à esquerda e fechado à direita, com exceção do primeiro intervalo, [ 0 ; 10 ], fechado em ambos os extremos.
Para os gráficos de distribuição das notas das questões discursivas, foram consideradas mais duas categorias : questão em branco21 e nota zero. 20 Uma definição dessas estatísticas pode ser encontrada no Glossário. 21 Nesse grupo estão incluídas também as respostas classificadas como nulas ou desconsideradas. 6.1 ESTATÍSTICAS BÁSICAS DA PROVA Esta seção apresenta estatísticas selecionadas e histogramas da nota geral ( 6.1.1 ) e de cada componente : Formação Geral ( 6.1.2 ) e Conhecimento Específico ( 6.1.3 ).
São também apresentadas estatísticas selecionadas de subpopulações caracterizadas por Grande Região, Categoria Administrativa, Organização Acadêmica e Modalidade de Ensino. 6.1.1 Estatísticas Básicas Gerais A Tabela 6.1 apresenta as Estatísticas Básicas da prova, por Grande Região segundo Modalidade de Ensino e Habilitação, dos estudantes concluintes de Química ( Licenciatura/Bacharelado ).
As análises serão feitas primeiramente para os cursos de Licenciatura e depois para os de Bacharelado.
Cada habilitação é desagregada internamente em cursos Presenciais e cursos a Distância.
A Média das notas da prova como um todo dos alunos de Licenciatura ( nas seções seguintes serão analisados os Componentes de Formação Geral e de Conhecimento Específico ), foi 38,1, sendo que os alunos da região Norte obtiveram a Média mais baixa ( 33,0 ), e os da região Sul obtiveram a Média mais alta ( 42,0 ).
As demais médias foram : 36,1 na região Nordeste, 40,9 na região Sudeste e 36,8 na região Centro-Oeste.
Já a Média das notas para os alunos de cursos presenciais de Licenciatura foi 38,2 sendo que os alunos da região Norte obtiveram a Média mais baixa ( 33,2 ), e os da região Sul obtiveram a Média mais alta ( 42,0 ).
As demais médias foram : 36,0 na região Nordeste, 41,8 na região Sudeste e 36,9 na região Centro-Oeste.
O Desvio padrão para o Brasil como um todo foi 12,0, sendo o maior Desvio padrão encontrado na região Sul ( 12,3 ) e o menor, na região Norte ( 10,5 ), indicando uma dispersão um pouco menor das notas dessa região.
A região que obteve a maior nota Máxima foi a Sudeste ( 86,8 ), ao passo que a região que atingiu a menor nota Máxima foi a Norte ( 65,7 ).
A Mediana do Brasil como um todo foi 37,7, sendo a maior Mediana obtida na região Sudeste ( 42,1 ), e a menor obtida na Norte ( 32,0 ).
A nota Mínima foi zero em todas as regiões, com exceção das regiões Sudeste ( 9,9 ) e Centro-Oeste ( 9,2 ).
Considerando -se as notas segundo Grande Região, observa -se que existe diferença estatisticamente significativa ao nível de 95 \% entre as maiores Médias, obtida nas regiões Sudeste ( 41,8 ) e Sul ( 42,0 ), e as médias de todas as demais regiões²². ²² Todas as comparações deste capítulo utilizam os intervalos de 95 \%.
Os erros-padrão da média que possibilitam os testes estão disponíveis nas tabelas.
Já a Média das notas da prova dos alunos de cursos a Distância de Licenciatura foi 36,8, mais baixa do que a dos alunos de cursos presenciais, sendo que os alunos da região Centro-Oeste obtiveram a Média mais baixa ( 30,4 ), e os da região Sudeste obtiveram a Média mais alta ( 37,4 ).
Não foram oferecidos cursos a Distância na região Sul.
As demais médias foram : 30,6 na região Norte e 36,7 na região Nordeste.
O Desvio padrão para o Brasil como um todo foi 12,6.
A região que obteve a maior nota Máxima foi a Nordeste ( 81,1 ), ao passo que a região que atingiu a menor nota Máxima foi a Centro-Oeste ( 30,4 ).
A Mediana do Brasil como um todo foi 35,3, sendo a maior Mediana obtida na região Sudeste ( 36,2 ), e a menor obtida na região Norte ( 30,1 ).
A nota Mínima foi zero na região Sudeste.
Considerando -se as notas segundo Grande Região, observa -se que existe diferença estatisticamente significativa ao nível de 95 \% entre a segunda menor Média, obtida na região Norte ( 30,6 ), e a maior Média, obtida na região Sudeste ( 37,4 ).
Cumpre notar que ainda que as notas dos cursos presenciais sejam comparáveis com as dos cursos a Distância dentro de uma mesma habilitação, não é possível comparar as notas obtidas pelos alunos de Licenciatura e Bacharelado, já que as provas são, em parte, diferentes.
A Média das notas da prova como um todo para os alunos de cursos de Bacharelado, foi 37,9, sendo que os alunos da região Norte obtiveram a Média mais baixa ( 31,1 ), e os da região Sudeste obtiveram a Média mais alta ( 38,3 ).
As demais médias foram : 37,4 na região Nordeste, 37,9 na região Sul e 36,6 na região Centro-Oeste.
A região que obteve a maior nota Máxima entre os alunos dos cursos de Bacharelado foi a Sudeste ( 86,9 ) e a com menor nota Máxima foi a Norte ( 56,0 ).
A Mediana do Brasil como um todo foi 36,7, sendo a maior Mediana obtida na região Sudeste ( 37,0 ), e a menor obtida na Norte ( 32,1 ).
A nota Mínima foi zero em três das regiões.
As exceções foram a região Norte com 11,1 e a região Nordeste com 7,5.
Considerando -se as notas segundo Grande Região, observa -se que existe diferença estatisticamente significativa ao nível de 95 \% entre a menor Média, obtida na região Norte ( 31,1 ) e a Média das demais regiões.
Não há registro de participantes na Modalidade Educação a Distância de Bacharelado.
Dessa forma, os valores correspondentes à Educação Presencial coincidem com os valores identificados como totais.
O comportamento das notas dos estudantes de todo o Brasil pode ser observado no Gráfico 6.1 que apresenta um histograma com a distribuição das mesmas segundo a habilitação.
Essa é uma distribuição unimodal para ambas as habilitações, com as modas no intervalo ( 30 ; 40 ] para os alunos de cursos de Licenciatura e de cursos de Bacharelado.
A Tabela 6.2 apresenta informações referentes às Estatísticas Básicas da prova desagregadas por Categoria Administrativa e por Organização Acadêmica segundo a Habilitação.
Os estudantes de Licenciatura de IES Públicas obtiveram Média mais alta ( 38,3 ), e os das IES Privadas obtiveram Média mais baixa ( 37,1 ) que a Média nacional para esta habilitação ( 38,1 ).
Observa -se que não existe diferença estatisticamente significativa entre as médias das notas das IES Públicas e Privadas.
O Desvio padrão para as IES Privadas ( 11,8 ) foi inferior ao do Brasil como um todo ( 12,0 ), indicando uma dispersão menor das notas nesta Categoria Administrativa do que para o conjunto de alunos de cursos de Licenciatura.
No tocante à Organização Acadêmica, alunos de Licenciatura de Universidades obtiveram nota média mais alta ( 39,0 ) do que os de Faculdades ( 36,5 ), que, por sua vez, obtiveram nota média mais alta do que os dos CEFET/IFET ( 36,4 ) e de Centros Universitários ( 33,8 ).
Existe diferença estatisticamente significativa entre a maior Média, obtida nas Universidades, e as médias das demais organizações acadêmicas.
Já em relação à Categoria Administrativa, os estudantes de Bacharelado das IES Públicas obtiveram Média mais alta ( 40,6 ), e os das IES Privadas obtiveram Média mais baixa ( 32,2 ) que a Média nacional para esta habilitação ( 37,9 ).
A diferença entre a Média de IES Públicas e IES Privadas é estatisticamente significativa.
O Desvio padrão para as IES Públicas ( 12,5 ) foi quase o mesmo que o do Brasil como um todo ( 12,4 ), e para IES Privadas ( 9,8 ) foi inferior ao do Brasil, indicando uma dispersão menor das notas nestas Categorias Administrativas do que para o conjunto de alunos de cursos de Bacharelado.
No tocante à Organização Acadêmica, alunos de CEFET/IFET obtiveram nota mais alta ( 39,4 ) do que os de Universidades ( 38,6 ), os de Faculdades ( 31,9 ) e do que os de Centros Universitários ( 31,6 ).
Observa -se que existe diferença estatisticamente significativa das maiores médias de Universidades e CEFET/IFET em relação às menores médias de Centros Universitários e Faculdades. 6.1.2 Estatísticas Básicas no Componente de Formação Geral A Tabela 6.3 apresenta as Estatísticas Básicas em relação ao componente da prova que avalia a Formação Geral dos estudantes concluintes por Grande Região, segundo Habilitação e Modalidade de ensino.
Os alunos de cursos de Licenciatura obtiveram desempenho médio de 52,2.
Quanto à variabilidade, o Desvio padrão das notas dos estudantes do Brasil, como um todo, foi 16,8.
A maior Média foi obtida na região Sul ( 57,4 ), e a menor, na região Norte ( 47,7 ).
As demais médias foram : 50,2 na região Nordeste, 54,4 na região Sudeste e 51,1 na região Centro-Oeste.
Já os alunos de cursos presenciais de Licenciatura obtiveram desempenho médio de 52,3.
Quanto à variabilidade, o Desvio padrão das notas dos estudantes do Brasil, como um todo, foi 16,7.
A maior Média foi obtida na região Sul ( 57,4 ), e a menor, na região Norte ( 47,8 ).
As demais médias foram : 50,1 na região Nordeste, 55,2 na região Sudeste e 51,1 na região Centro-Oeste.
Já o maior Desvio padrão foi obtido na região Sul ( 17,0 ), e o menor, nas regiões Norte e Centro-Oeste ( 16,1 em ambas ).
Os demais desvios padrões foram iguais a 16,5 para as regiões Nordeste e Sudeste.
A maior nota no Componente de Formação Geral da prova de alunos de cursos presenciais de Licenciatura foi 97,8, obtida por, pelo menos, um aluno na região Sudeste.
A menor nota Máxima foi obtida na região Norte ( 86,3 ).
A Mediana do Brasil, como um todo, foi 53,3, sendo a menor Mediana encontrada na região Norte ( 49,1 ), e a maior encontrada na região Sul ( 59,2 ).
A nota Mínima nesta parte foi zero em todas as regiões, com exceção da região Centro-Oeste com 7,5.
Considerando -se as notas segundo Grande Região, observa -se que existe diferença estatisticamente significativa entre a menor Média das notas do Componente de Formação Geral, obtida na região Norte ( 47,8 ), e as médias das demais regiões.
Já a Média das notas do Componente de Formação Geral dos alunos de cursos de Licenciatura a Distância foi 51,4, mais baixa do que a dos alunos de cursos presenciais, sendo que os alunos da região Centro-Oeste obtiveram a Média mais baixa ( 39,2 ), e os da região Nordeste obtiveram a Média mais alta ( 53,3 ).
Como já comentado, não foram oferecidos cursos a Distância na região Sul.
As demais médias foram : 46,9 na região Norte e 51,5 na região Sudeste.
O Desvio padrão para o Brasil como um todo foi 17,6, sendo o maior Desvio padrão encontrado na região Sudeste ( 18,2 ) e o menor, na região Norte ( 14,9 ).
A região que obteve a maior nota Máxima foi a Nordeste ( 92,6 ), ao passo que a região que atingiu a menor nota Máxima foi a Centro-Oeste ( 39,2 ).
A Mediana do Brasil como um todo foi 52,5, sendo a maior Mediana obtida na região Sudeste ( 53,4 ) e a menor obtida na Centro-Oeste ( 39,2 ).
A nota Mínima foi zero na região Sudeste, 11,1 na Norte, 39,2 na região Centro-Oeste e 18,1 na região Nordeste.
A Média das notas dos alunos de cursos de Bacharelado, foi 56,5, sendo que os alunos da região Norte obtiveram a Média mais baixa ( 51,8 ), e os da região Sudeste obtiveram a Média mais alta ( 57,3 ).
As demais médias foram : 54,3 na região Nordeste, 56,2 na região Sul e 55,2 na região Centro-Oeste.
A região que obteve a maior nota Máxima no curso de Bacharelado foi a Sudeste ( 97,0 ) e a com menor nota Máxima foi a Norte ( 86,5 ).
A Mediana do Brasil como um todo foi 57,6, sendo a maior Mediana obtida na região Sudeste ( 58,6 ), e a menor obtida na Norte ( 49,7 ).
A nota Mínima foi zero em três das regiões.
As exceções foram a região Norte com 15,0 e a região Nordeste com 13,3.
Considerando -se as notas segundo Grande Região, observa -se que existe diferença estatisticamente significativa ao nível de 95 \% entre a menor das Médias, obtida na região Norte, e a maior delas, a da região Sudeste.
Não há registro de participantes de cursos a Distância de Bacharelado.
O Gráfico 6.2 propicia a avaliação do desempenho dos estudantes do Componente de Formação Geral, a partir do histograma da distribuição das notas correspondentes.
A distribuição é unimodal, com moda em ( 50 ; 60 ] para ambas as Habilitações.
No Bacharelado, o intervalo modal é seguido de perto pelo intervalo ( 60, 70 ].
Este intervalo modal está mais à direita do que o intervalo modal da distribuição de notas da prova.
Nota -se, ainda, que no Gráfico 6.2, as notas apresentam dispersão maior, quando comparadas a do Gráfico 6.1 ( distribuição das notas da prova ), confirmado pela comparação dos desvios padrões.
Na Tabela 6.4 são apresentadas as informações referentes ao desempenho dos concluintes do Componente de Formação Geral, em diferentes agregações : Categoria Administrativa, Organização Acadêmica e Habilitação.
Primeiramente, com respeito aos alunos de Licenciatura e Categorias Administrativas, a maior Média foi obtida por estudantes de IES Públicas ( 52,4 ) e a menor, pelos de IES Privadas ( 51,8 ).
A diferenças entre essas médias não é estatisticamente signicifcativa.
A ordenação por Organização Acadêmica foi : Universidades ( 52,9 ), Faculdades ( 51,5 ), CEFET/IFET ( 51,2 ) e Centros Universitários ( 48,8 ).
A diferença entre a maior Média das Universidades e as duas menores, de Centros Universitários e CEFET/IFET, é estatisticamente significativa.
Entre os alunos de Bacharelado, a situação das médias com relação à Categoria Administrativa se mantém ( maior para IES Públicas, 57,9, menor para IES Privadas, 53,6 ), mas a diferença entre elas é estatisticamente significativa.
No Bacharelado, a maior Média foi obtida nos CEFET/IFET ( 61,5 ), seguidos, sucessivamente, pelas Universidades ( 56,8 ), pelos Centros Universitários ( 53,5 ), e por fim, pelas Faculdades ( 53,3 ).
A diferença da maior Média, obtida em CEFET/IFET e as demais médias é estatisticamente significativa. 6.1.3 Estatísticas Básicas do Componente de Conhecimento Específico A Tabela 6.5 apresenta as Estatísticas Básicas referentes ao Componente de Conhecimento Específico da área de Química ( Licenciatura/Bacharelado ) por Grande Região, segundo Habilitação e Modalidade de Ensino.
Quanto ao desempenho de Conhecimento Específico dos alunos dos cursos de Licenciatura, a Média foi 33,3.
A maior Média foi obtida na região Sul ( 36,8 ), e a menor, na região Norte ( 28,1 ).
As demais médias foram : 31,3 na região Nordeste, 36,4 na região Sudeste e 32,1 na região Centro-Oeste.
Para os alunos dos cursos presenciais, a Média do desempenho de Conhecimento Específico foi semelhante, 33,4.
A maior Média foi obtida na região Sudeste ( 37,4 ), e a menor, na região Norte ( 28,3 ).
As demais médias foram : 31,3 na região Nordeste, 36,8 na região Sul e 32,1 na região Centro-Oeste.
Quanto à variabilidade das notas, o Desvio padrão do Brasil, como um todo, foi 12,8, sendo o maior Desvio padrão observado nas regiões Sudeste e Sul ( 13,0 em ambas ), e o menor, na região Norte ( 11,2 ).
Os demais desvios foram : 12,4 na região Nordeste e 11,8 na região Centro-Oeste.
A Mediana das notas dos estudantes de cursos presenciais de Licenciatura de todo o Brasil foi 32,2.
A maior Mediana ocorreu na região Sudeste ( 36,9 ), e a menor, na região Norte ( 27,0 ).
As demais medianas foram : 30,9 nas regiões Nordeste e Centro-Oeste e 35,8 na região Sul.
A nota Máxima do Brasil, como um todo, foi 89,2, sendo obtida por, pelo menos, um aluno na região Sudeste.
As demais notas máximas foram : 69,0 na região Norte, 76,5 na região Nordeste, 83,7 na região Sul e 71,7 na região Centro-Oeste.
A nota Mínima foi zero em todas as regiões, com exceção das regiões Sudeste e Centro-Oeste, ambas com 3,8.
Observa -se que existe diferença estatisticamente significativa entre a menor Média das notas do Componente de Conhecimento Específico, encontrada na região Norte, em relação às notas das demais regiões.
Já a Média do desempenho de Conhecimento Específico dos alunos dos cursos a Distância de Licenciatura foi um pouco menor, 31,8, quando comparado com o Presencial, que foi 33,4.
A maior Média foi obtida na região Sudeste ( 32,6 ), e a menor, na região Norte ( 25,1 ).
As demais médias foram : 31,2 na região Nordeste e 27,5 na Centro-Oeste.
Não há registro de participantes no Ensino a Distância de Licenciatura na região Sul.
O Desvio padrão de todos os alunos foi 13,7, sendo o maior Desvio padrão observado na região Nordeste ( 16,4 ), e o menor, na região Norte ( 9,8 ).
A Mediana das notas dos estudantes de cursos a Distância de Licenciatura foi 30,9, mesma Mediana da região Sudeste, a maior Mediana entre as regiões.
A menor Mediana foi na região Norte ( 23,7 ).
As demais medianas foram : 28,8 na região Nordeste e 27,5 na região Sul.
A nota Máxima do Brasil, como um todo, foi 83,2, sendo obtida por, pelo menos, um aluno na região Nordeste.
As demais notas máximas foram : 46,3 na região Norte, 72,8 na região Sudeste e 27,5 na região Centro-Oeste.
A nota Mínima foi zero, o mesmo valor em duas regiões : Nordeste e Sudeste.
A nota Mínima da região Norte foi 11,6 e a da Centro-Oeste, 27,5.
Como já foi comentado, não há registro de participante nessa Modalidade e Habilitação para a região Sul.
Para os alunos de cursos de Bacharelado, a Média do desempenho de Conhecimento foi 31,7.
A maior Média foi obtida na região Sudeste ( 32,0 ), e a menor, na região Norte ( 24,1 ).
As demais médias foram : 31,8 nas regiões Nordeste e Sul, e 30,4 na região Centro-Oeste.
O Desvio padrão de todos os alunos de Bacharelado foi 13,7, sendo o maior Desvio padrão observado na região Sudeste ( 14,1 ), e o menor, na região Norte ( 11,0 ).
Os demais desvios foram : 13,5 na região Nordeste, 13,6 na região Sul e 11,5 na região Centro-Oeste.
A Mediana das notas dos estudantes de Bacharelado foi 29,8.
A maior Mediana ocorreu na região Sul ( 30,3 ), e a menor, na região Norte ( 24,0 ).
As demais medianas foram : 30,1 na região Nordeste, 29,8 nas regiões Sudeste e Centro-Oeste.
A nota Máxima do Brasil, como um todo, foi 86,9, obtida por pelo menos, um aluno na região Sudeste.
As demais notas máximas foram : 50,8 na região Norte, 74,8 na região Nordeste, 79,9 na região Sul e 62,9 na região Centro-Oeste.
A nota Mínima foi 0,0 em três regiões : Sudeste, Sul e Centro-Oeste.
Na região Norte, a Mínima foi 4,1, e na Nordeste foi 5,6.
Não há registro de cursos na Modalidade Ensino a Distância na Habilitação Bacharelado.
Sendo assim, todos os registros numéricos totais coincidem com os de Educação Presencial.
Observa -se que para os alunos de cursos de Bacharelado, existe diferenças estatisticamente significativas da menor Média das notas do Componente de Conhecimento Específico, obtida na região Norte, em relação às demais regiões.
Assim como os Gráficos 6.1 e 6.2, o Gráfico 6.3, apresentado a seguir, permite uma avaliação do desempenho de concluintes de Química ( Licenciatura/Bacharelado ) segundo a Habilitação, em relação ao Componente de Conhecimento Específico com um histograma da distribuição das notas correspondentes.
Esta também é uma distribuição unimodal para ambas as habilitações, e os grupos modais são, respectivamente, ( 30 ; 40 ] e ( 20 ; 30 ], para estudantes de Licenciatura e Bacharelado. ambas as habilitações, e os grupos modais são, respectivamente, ( 30 ; 40 ] e ( 20 ; 30 ], para estudantes de Licenciatura e Bacharelado.
A Tabela 6.6 apresenta uma comparação dos resultados em relação à Categoria Administrativa e à Organização Acadêmica segundo a Habilitação, agora levando em conta o desempenho de alunos do Componente de Conhecimento Específico da prova.
No que se refere à Organização Acadêmica, para os alunos de Licenciatura, a maior Média foi a de Universidades ( 34,3 ), vindo a seguir a de Faculdades e a dos CEFET/IFET ( 31,5 cada uma dessas organizações ) e, por último, a dos Centros Universitários ( 28,8 ).
O maior Desvio padrão, e acima do valor para o Brasil como um todo ( 13,0 ), foi o das Universidades ( 13,4 ), organizações que também obtiveram a maior nota Máxima, igual a 89,2, com os CEFET/IFET em segundo lugar, com nota Máxima 83,5.
As Faculdades obtiveram nota Máxima igual a 69,7 e os Centros Universitários, 54,1.
As medianas foram : 33,4 nas Universidades, a maior Mediana, 27,8 nos Centros Universitários, a menor delas, 30,9 em Faculdades, e CEFET/IFET.
A nota Mínima foi zero para Universidades e CEFET/IFET e 3,8 para Centros Universitários e Faculdades.
Quanto à Categoria Administrativa, observa -se um comportamento semelhante àquele da parte de Formação Geral e à prova como um todo, ou seja, a Média das IES Públicas ( 33,6 ) é maior do que a Média das IES Privadas ( 32,1 ), e a diferença entre elas é estatisticamente significativa.
Para os alunos de Bacharelado, a maior Média, considerando -se a Organização Acadêmica, foi de Universidades ( 32,6 ), e a menor, a dos Centros Universitários ( 24,4 ).
As Universidades obtiveram a maior nota Máxima, igual a 86,9.
As medianas foram : 30,8 nas Universidades, 24,0 nos Centros Universitários, 22,8 nas Faculdades e 30,3 nos CEFET/IFET.
A nota Mínima foi zero nas Universidades e nas Faculdades, foi 4,1 nos Centros Universitários e foi 5,8 nos CEFET/IFET.
Quanto à Categoria Administrativa, observa -se um comportamento semelhante àquele da parte de Formação Geral e à prova como um todo para os alunos de Bacharelado, ou seja, a Média das IES Públicas ( 34,8 ) é maior do que a das IES Privadas ( 25,0 ), com diferença estatisticamente significativa entre essas médias. 6.2 ANÁLISE DAS QUESTÕES OBJETIVAS Esta seção apresenta estatísticas selecionadas e histogramas das Questões Objetivas de Formação Geral para Licenciatura ( 6.2.1 ) e para Bacharelado ( 6.2.2 ), e das Questões Objetivas de Conhecimento Específico para Licenciatura ( 6.2.3 ) e para Bacharelado ( 6.2.4 ).
São também apresentadas e comparadas as médias das subpopulações caracterizadas por Grande Região. 6.2.1 Componente de Formação Geral – Química ( Licenciatura ) A Tabela 6.7a apresenta as Estatísticas Básicas relativas às oito questões objetivas do componente da prova que abrange a Formação Geral dos estudantes de Química ( Licenciatura ) segundo a Modalidade de ensino.
A Média dos alunos nestas questões foi 50,0.
A menor Média foi encontrada na região Norte ( 44,5 ), e a maior, na região Sul ( 56,3 ).
As demais médias foram : 47,8 na região Nordeste, 52,5 na região Sudeste e 48,0 na região Centro-Oeste.
Para os alunos de cursos presenciais, a Média foi 50,1.
A menor Média foi encontrada na região Norte ( 44,5 ), e a maior, na região Sul ( 56,3 ).
As demais médias foram : 47,7 na região Nordeste, 53,4 na região Sudeste e 48,1 na região Centro-Oeste.
O Desvio padrão do Brasil foi 20,3, sendo o maior Desvio padrão encontrado na região Sul ( 20,7 ), e o menor, na região Nordeste ( 19,6 ).
Os demais desvios foram : 20,5 na região Norte, 20,0 na região Sudeste e 19,7 na região Centro-Oeste.
As medianas do Brasil como um todo e de todas as cinco regiões foi 50,0.
A nota Máxima 100,0 foi alcançada em todas as grandes Regiões.
As notas mínimas ( 0,0 ) foram iguais para todas as regiões.
Já para os alunos de cursos a Distância, a Média foi menor, 48,8.
A menor Média foi encontrada na região Centro-Oeste ( 25,0 ), e a maior, na região Nordeste ( 50,2 ).
As demais médias foram : 43,8 na região Norte e 49,1 na região Sudeste.
Não há registro de participantes de Educação a Distância na região Sul.
O Desvio padrão do Brasil para esta modalidade de ensino nos cursos de Licenciatura foi 20,5, sendo o maior Desvio padrão encontrado na região Sudeste ( 20,6 ), e o menor, na região Norte ( 17,9 ).
As medianas do Brasil como um todo e para duas das regiões ( Nordeste e Sudeste ) foram 50,0.
Para as outras duas regiões foi mais baixa, 43,8 para o Norte e 25,0 para o Centro-Oeste.
A nota Máxima 100,0 foi alcançada no Nordeste e no Sudeste.
No Norte, a nota Máxima foi 75,0 e no Centro-Oeste foi 25,0.
A nota Mínima ( 0,0 ) foi registrada na região Sudeste.
Nas regiões Norte e Nordeste, tal nota foi 12,5 e, na região Centro-Oeste, foi 25,0.
Como já foi mencionado, não houve registro de Educação a Distância na região Sul.
A Tabela 6.8a apresenta o Índice de Facilidade e o Índice de Discriminação ( ponto-bisserial ) para cada uma das questões objetivas do Componente de Formação Geral.
Quanto ao Índice de Facilidade, foram usadas as seguintes cores para diferenciar o nível de dificuldade da questão :  Azul para as questões classificadas com índice Muito fácil ( \&gt ; =0,86 ), verde para as questões classificadas com índice Fácil ( 0,61 a 0,85 ), amarelo para as questões classificadas com Médio ( 0,41 a 0,60 ), vermelho para as questões classificadas com Difícil ( 0,16 a 0,40 ) e roxo para as questões classificadas com Muito difícil ( \&lt ; =0,15 ).
Já quanto ao Índice de Discriminação, foram usadas as seguintes cores para qualificar a questão :  As questões classificadas com índice Fraco receberam a cor vermelho ( \&lt ; =0,19 ), as classificadas com Médio receberam a cor amarelo ( 0,20 a 0,29 ), as classificadas com Bom receberam a cor verde ( 0,30 a 0,39 ) e as classificadas com Muito bom ( \&gt ; =0,40 ) receberam a cor azul.
As questões objetivas do Componente de Formação Geral, segundo o Índice de Facilidade, foram assim avaliadas : das oito questões, nenhuma teve o Índice de Facilidade classificado como Muito fácil e duas questões foram tidas como Fácil.
Três questões foram consideradas com índice de dificuldade Médio, situando -se no intervalo entre 0,41 e 0,60 do Índice de Facilidade, ou seja, houve entre 41,0 \% e 60,0 \% de acertos, enquanto outras três questões foram classificadas na categoria Difícil, situando -se no intervalo entre 0,16 e 0,40.
Nenhuma questão apresentou menos de 15 \% de acertos, razão pela qual não houve questão classificada como Muito difícil.
O Índice de Facilidade variou de 0,20 a 0,84, e o de Discriminação, de 0,39 a 0,49.
A Tabela 6.9a apresenta a distribuição das questões considerando -se simultaneamente a classificação dos dois índices.
As sete questões com Índice de Discriminação Muito bom figuraram entre três níveis de dificuldade Fácil, Médio e Difícil : duas classificadas na categoria Fácil ( questões 4 e 6 ) do Índice de Facilidade, duas na categoria Médio ( questões 5 e 8 ) e três na categoria Difícil ( questões 1, 3 e 7 ).
Em particular, a questão 4 foi a que apresentou o maior poder discriminatório, com índice 0,49, e foi considerada Fácil, com uma proporção de 0,67 acertos.
O máximo de acertos foi alcançado pela questão 6 com um Índice de Facilidade de 0,84.
A questão de número 1 apresentou um Índice de Facilidade de 0,20 ( Difícil ), ou seja, um quinto dos estudantes conseguiu resolvê -la, dentro do universo de participantes.
Seu Índice de Discriminação foi Muito bom ( 0,41 ).
O Gráfico 6.4a, para exemplificar, analisa o comportamento da questão de número 4 de Formação Geral.
Trata -se de uma questão considerada Fácil, com relação à facilidade e uma das que obteve o maior Índice de Discriminação dessa parte da prova ( 0,49 ).
Neste gráfico, cada uma das cinco curvas representa o percentual de respostas em determinada alternativa da questão, em função do número de acertos dos estudantes nessa parte da prova ( Formação Geral/Múltipla Escolha ), antes de possíveis eliminações pelo critério do ponto-bisserial.
A curva em verde corresponde à alternativa B, a alternativa correta para esta questão.
Observa -se que entre os estudantes com menor número de acertos, nessa parte do exame, a situação mais frequente foi a escolha da alternativa E ( em vermelho ), incorreta.
Por exemplo, entre os estudantes que acertaram três questões, 20,1 \% escolheram a alternativa E ( em vermelho ), mais da metade ( 55,1 \% ) escolheram a alternativa correta B ( em verde ), 3,6 \% escolheram a alternativa C ( preto ), 8,0 \%, a alternativa A ( azul ) e 13,1 \%, a D ( roxo ).
Entre os que acertaram três respostas entre as questões de múltipla escolha de Formação Geral, nenhum deixou a questão em branco e 0,1 \% marcou mais de uma alternativa, invalidando a questão.
À medida que o número de acertos aumenta, indicando desempenho melhor nessa parte da prova, aumenta concomitantemente a proporção de estudantes que selecionaram a alternativa correta B, atingindo 100 \% para os estudantes com 8 acertos.
Essa análise permite verificar como a questão discriminou os grupos de desempenho, justificando o alto índice obtido na questão.
Cumpre notar que não é possível inferir deste gráfico nem o índice de facilidade ( que seria uma média da proporção ponderada pela quantidade de alunos com cada uma das notas ), nem o índice de discriminação ponto-bisserial, por razão equivalente.
No caso extremo no qual a grande concentração dos acertos dos alunos fosse abaixo de 4, o índice de facilidade seria obrigatoriamente abaixo de 40 \% ( neste exemplo ).
Caso a concentração fosse em 6 acertos ou mais, o índice seria obrigatoriamente acima de 60 \%.
Os gráficos relativos às demais questões de Formação Geral constam do Anexo Ia. 6.2.2 Componente de Formação Geral – Química ( Bacharelado ) A Tabela 6.7b apresenta as Estatísticas Básicas relativas às oito questões objetivas do componente da prova que abrange a Formação Geral dos estudantes de Química ( Bacharelado ) segundo a Modalidade de ensino.
A Média dos alunos de Bacharelado nestas questões foi 56,4.
A menor Média foi encontrada na região Norte ( 50,4 ), e a maior, nas regiões Sudeste e Sul ( 57,0 em ambas ).
As demais médias foram : 53,8 na região Nordeste e 55,2 na região Centro-Oeste.
O Desvio padrão do Brasil foi 20,4, sendo o maior Desvio padrão encontrado na região Norte ( 21,2 ), e o menor, na região Centro-Oeste ( 19,4 ).
Os demais desvios foram : 20,1 na região Nordeste e 20,5 nas regiões Sudeste e Sul.
As medianas do Brasil como um todo e de três das cinco regiões ( Sudeste, Sul e Centro-Oeste ) foi 62,5.
As medianas das regiões Norte e Nordeste foram mais baixas, 50,0.
A nota Máxima 100,0 foi alcançada em todas as grandes Regiões, com exceção da região Norte cuja nota Máxima foi 87,5.
As notas mínimas ( 0,0 ) foram iguais para todas as regiões, sem exceção.
Como já foi mencionado, não houve registro de cursos a Distância de Bacharelado, levando a identificação dos registros totais com os registros da Educação Presencial.
A Tabela 6.8b apresenta o Índice de Facilidade e o Índice de Discriminação ( ponto-bisserial ) para cada uma das questões objetivas do Componente de Formação Geral.
A convenção das cores é a mesma da Tabela 6.8a.
As questões objetivas do Componente de Formação Geral, segundo o Índice de Facilidade, foram assim avaliadas : das oito questões, uma teve o Índice de Facilidade classificado como Muito fácil e duas questões foram tidas como Fácil.
Três questões foram consideradas com índice de dificuldade Médio, enquanto duas questões foram classificadas na categoria Difícil, situando -se no intervalo entre 0,16 e 0,40.
Nenhuma das questões apresentou menos de 15 \% de acertos, portanto, não houve questão classificada como Muito difícil.
O Índice de Facilidade variou de 0,26 a 0,90, e o de Discriminação, de 0,37 a 0,49.
A Tabela 6.9b apresenta a distribuição das questões considerando -se simultaneamente a classificação dos dois índices, semelhantemente ao apresentado na Tabela 6.9a, mas aqui para os alunos de Bacharelado.
As sete questões com Índice de Discriminação Muito bom figuraram entre quatro níveis de dificuldade Muito fácil, Fácil, Médio e Difícil : uma classificada na categoria Muito fácil ( a questão 6 ) do Índice de Facilidade, duas, na categoria Fácil ( questões 4 e 5 ), duas, na categoria Médio ( questões 3 e 8 ) e, por fim, duas, na categoria Difícil ( questões 1 e 7 ).
Em particular, a questão 8 foi uma das que apresentou o maior poder discriminatório, com índice 0,49, e foi considerada Médio em, com uma proporção de 0,52 acertos.
O máximo de acertos foi alcançado pela questão 6 com um Índice de Facilidade de 0,90.
O Gráfico 6.4b, para exemplificar, analisa o comportamento da questão de número 4 de Formação Geral.
Trata -se de uma questão considerada Fácil, com relação à facilidade e a com Índice de Discriminação Muito bom.
Semelhantemente ao apresentado no Gráfico 6.4a, neste gráfico também, cada uma das cinco curvas representa o percentual de respostas em determinada alternativa da questão, em função do número de acertos dos estudantes nessa parte da prova ( Formação Geral/Múltipla Escolha ), antes de possíveis eliminações pelo critério do ponto-bisserial.
A curva em verde corresponde à alternativa B, a alternativa correta para esta questão.
Observa -se que entre os estudantes com menor número de acertos, nessa parte do exame, a situação mais frequente foi a escolha da alternativa E ( em vermelho ), incorreta.
Por exemplo, entre os estudantes que acertaram três questões, 16,1 \% escolheram a alternativa E ( rem vermelho ), mais da metade ( 55,9 \% ) escolheram a alternativa correta B ( em verde ), 14,5 \% escolheram a alternativa D ( roxo ), 7,7 \% escolheram a alternativa A ( em azul ) e 5,8 \% escolheram a alternativa C ( preto ).
Entre os que acertaram três respostas entre as questões de múltipla escolha de Formação Geral, nenhum deixou a questão em branco ou marcou mais de uma alternativa.
À medida que o número de acertos aumenta, indicando desempenho melhor nessa parte da prova, aumenta concomitantemente a proporção de estudantes que selecionaram a alternativa correta B, atingindo 100 \% para os estudantes com 8 acertos.
Essa análise permite verificar como a questão discriminou os grupos de desempenho, justificando o alto índice obtido na questão.
Os gráficos relativos às demais questões de Formação Geral constam do Anexo Ib. 6.2.3 Componente de Conhecimento Específico – Química ( Licenciatura ) A Tabela 6.10a apresenta as Estatísticas Básicas em relação às questões objetivas do Componente de Conhecimento Específico da prova de Química ( Licenciatura ) por Grande Região segundo a modalidade de ensino.
Nessa parte da prova, nenhuma das questões foi anulada pela Comissão Assessora da Área.
Como será discutido ainda nessa seção, cinco questões objetivas do Componente de Conhecimentos Específicos não foram usadas no cômputo das notas por terem sido descartadas pelo critério do ponto-bisserial.
Assim, as notas foram calculadas com base em 22 das 27 questões objetivas de Conhecimentos Específicos.
A Média deste componente para os alunos de cursos de Licenciatura foi de 37,3.
A menor Média foi observada na região Norte ( 32,0 ), e a maior, na região Sul ( 41,0 ).
O Desvio padrão de todo o Brasil foi 14,1, sendo o menor Desvio padrão encontrado na região Norte ( 12,4 ), e o maior, na região Sudeste ( 14,2 ).
Na região Nordeste, o Desvio padrão foi 13,7, na região Sul foi 14,1 e na Centro-Oeste, 13,0.
A situação é um pouco diferente quando desagregamos por Modalidade de Ensino : a Média deste componente para os alunos de cursos presenciais de Licenciatura foi de 37,4.
A menor Média foi observada na região Norte ( 32,2 ), e a maior, na região Sudeste ( 41,4 ).
O Desvio padrão de todo o Brasil foi 14,0, sendo o menor Desvio padrão encontrado na região Norte ( 12,5 ), e o maior, na região Sul ( 14,1 ).
A Mediana de todo o Brasil foi 36,4, a mesma das regiões Nordeste e Centro-Oeste.
Nas regiões Sudeste e Sul, a Mediana foi a mesma ( 40,9 ) e maior que a nacional.
Na região Norte ( 31,8 ) a Mediana foi menor.
A nota Máxima da prova foi 95,5, obtida nas questões objetivas do Componente de Conhecimento Específico, por, pelo menos, um aluno da região Sudeste.
No Nordeste a nota Máxima foi 81,8, na região Sul foi 86,4 e na região Centro-Oeste, 77,3.
A nota Mínima foi zero nas regiões Norte, Nordeste e Sul, e foi 4,5 nas regiões Sudeste e Centro-Oeste.
Já para os alunos de cursos a Distância de Licenciatura, a Média deste componente foi mais baixa, 36,0.
A menor Média foi observada na região Norte ( 29,1 ), e a maior, na região Sudeste ( 36,8 ).
A Média da região Nordeste foi 35,4 e da região Centro-Oeste foi 31,8.
Não houve registro de participantes na região Sul na modalidade Educação a Distância da Licenciatura.
O Desvio padrão de todo o Brasil foi 14,8, sendo o menor Desvio padrão encontrado na região Norte ( 11,0 ), e o maior, na região Nordeste ( 18,0 ).
A Mediana de todo o Brasil foi 36,4, a mesma da região Sudeste.
A nota Máxima da prova foi 86,4, também a mesma obtida na região Nordeste.
Na região Norte, a nota Máxima foi 50,0, na região Sudeste foi 77,3 e na Centro-Oeste foi 31,8.
A nota Mínima ( 0,0 ) foi obtida nas regiões Nordeste e Sudeste.
A Tabela 6.11a apresenta os Índices de Facilidade e Discriminação ( ponto-bisserial ) das questões objetivas do Componente de Conhecimento Específico para os estudantes de Química ( Licenciatura ).
Para facilitar a diferenciação das questões, usam -se as mesmas cores da Tabela 6.8ae 6.8b para as diferentes classificações dos Índices de Facilidade e de Discriminação.
Observa -se que nenhuma das questões foi anulada pela CAA, e as análises dizem respeito às 27 questões válidas.
A partir do Índice de Facilidade obtido, pode -se concluir que mais do que a metade das questões objetivas da prova foi considerada, pelo menos, Difícil : das 27 questões, 16 foram classificadas como Difícil ou como Muito difícil.
Uma questão foi classificada como Fácil, e outras 10 consideradas Médio.
Nenhuma questão foi classificada como Muito fácil.
Já quanto aos índices de discriminação das questões objetivas do Componente de Conhecimento Específico da prova, tem -se como resultado a seguinte classificação : oito das 27 questões foram consideradas boas, enquanto duas delas tiveram Índice de Discriminação Muito bom.
Assim, para um pouco mais de um terço das questões – 10 em 27 – os Índices de Discriminação foram Bom ou Muito bom.
Dentre as demais, 12 delas foram classificadas como Médio e outras cinco como Fraco, sendo 17, por conseguinte, a quantidade de questões nos dois patamares mais baixos de discriminação.
Constata -se, assim, que a prova – no que se refere ao Componente de Conhecimento Específico – possuía média capacidade de discriminar entre aqueles que dominavam ou não o conteúdo.
A Tabela 6.12a apresenta a distribuição das questões considerando -se simultaneamente a classificação dos dois índices.
Dentre as questões que alcançaram os maiores índices de discriminação, classificadas com Índice de Discriminação Muito bom, uma delas, a de número 23, a com índice mais elevado ( 0,44 ), foi classificada como Médio quanto ao Índice de Facilidade, com 45 \% dos estudantes marcando a opção correta.
A questão de número 9 foi a mais difícil dentre as 27 questões específicas, com baixo Índice de Facilidade, apenas 13 \% de acertos.
Essa questão apresentou poder discriminatório médio, 0,20.
Destaca -se, também, a questão 21, com Índice de Facilidade 0,14, o que, em termos percentuais, corresponde a 14 \% de estudantes que responderam acertadamente.
Já 0,10 foi o seu de Índice de Discriminação.
Tal questão foi, portanto, pelo critério ponto-bisserial, considerada inadequada.
Por isso, a questão 21, além das questões 10, 13, 24 e 26, foram eliminadas do cômputo da nota final.
A Tabela 6.12a apresenta a distribuição das questões considerando -se simultaneamente a classificação dos dois índices.
Dentre as questões que alcançaram os maiores índices de discriminação, classificadas com Índice de Discriminação Muito bom, uma delas, a de número 23, a com índice mais elevado ( 0,44 ), foi classificada como Médio quanto ao Índice de Facilidade, com 45 \% dos estudantes marcando a opção correta.
A questão de número 9 foi a mais difícil dentre as 27 questões específicas, com baixo Índice de Facilidade, apenas 13 \% de acertos.
Essa questão apresentou poder discriminatório médio, 0,20.
Destaca -se, também, a questão 21, com Índice de Facilidade 0,14, o que, em termos percentuais, corresponde a 14 \% de estudantes que responderam acertadamente.
Já 0,10 foi o seu de Índice de Discriminação.
Tal questão foi, portanto, pelo critério ponto-bisserial, considerada inadequada.
Por isso, a questão 21, além das questões 10, 13, 24 e 26, foram eliminadas do cômputo da nota final.
A título de exemplo das análises do comportamento das questões objetivas, o Gráfico 6.5a analisa a questão 23 do Componente de Conhecimento Específico.
Esta questão foi considerada pelas respostas dos estudantes avaliados na prova como uma questão com nível de facilidade Médio, 0,45, ou seja, 45 \% dos estudantes assinalaram acertadamente a opção E, correspondente ao gabarito.
Como já comentado, seu Índice de Discriminação foi igual a 0,44, classificado como Muito bom, o maior valor de discriminação.
Neste gráfico também, cada uma das cinco curvas representa o percentual de respostas em determinada alternativa da questão 23, em função do número de acertos dos estudantes nessa parte da prova, antes de possíveis eliminações de questões pelo critério do ponto-bisserial.
A alternativa correta E, representada no gráfico pela curva em vermelho, foi escolhida em maiores proporções pelos alunos com desempenho melhor nessa parte da prova.
Já as alternativas incorretas, também denominadas distratores, foram selecionadas, principalmente, por aqueles com notas mais baixas.
Observa -se que a soma não é 100 \% por causa das questões não respondidas ou com mais de uma opção marcada.
Aqueles com nota zero, na sua quase totalidade deixaram esta questão em branco ou marcaram mais de uma alternativa, comportamento considerado inválido.
A proporção de alunos que selecionou a resposta correta E aumenta gradativamente, chegando a atingir 100 \% para 18, 19 e 21 acertos ou mais, enquanto a proporção dos que escolheram alternativas incorretas decai, a partir de um ou dois acertos, como função do número de acertos nessa parte da prova.
Os gráficos relativos às demais questões do Conhecimento Específico constam do Anexo I. 6.2.4 Componente de Conhecimento Específico – Química ( Bacharelado ) A Tabela 6.10b apresenta as Estatísticas Básicas em relação às questões objetivas do Componente de Conhecimento Específico da prova de Química ( Bacharelado ) por Grande Região segundo a Modalidade de Ensino.
Uma das questões foi anulada pela Comissão Assessora da Área e, como será discutido ainda nessa seção, cinco das questões objetivas do Componente de Conhecimentos Específicos não foram usadas no cômputo das notas por terem sido descartadas pelo critério do ponto-bisserial.
Assim, as notas foram calculadas com base em 21 das 27 questões objetivas de Conhecimentos Específicos.
Cabe destacar que não houve registro de participantes na Modalidade Educação a Distância.
Dessa forma, os registros feitos para Educação Presencial e os registros totais são coincidentes.
A Média deste componente para os alunos de cursos de Bacharelado foi de 33,3.
A menor Média foi observada na região Norte ( 25,4 ), e a maior, na região Sudeste ( 33,7 ).
O Desvio padrão de todo o Brasil foi 14,5, sendo o menor Desvio padrão encontrado na região Norte ( 12,0 ), e o maior, na região Sudeste ( 14,8 ).
A Mediana de todo o Brasil foi 33,3, a mesma nas regiões Nordeste, Sudeste, Sul e Centro-Oeste.
Na região Norte a Mediana foi menor, 23,8.
A nota Máxima da prova foi 95,2, obtida nas questões objetivas do Componente de Conhecimento Específico, por, pelo menos, um aluno da região Sudeste.
A nota Mínima foi zero nas regiões Sudeste, Sul e Centro-Oeste.
Nas regiões Norte e Nordeste, a nota Mínima foi 4,8.
A Tabela 6.11b apresenta os Índices de Facilidade e Discriminação ( ponto-bisserial ) das questões objetivas do Componente de Conhecimento Específico para os estudantes de Química ( Bacharelado ).
Para facilitar a diferenciação das questões, usaremos as mesmas convenções definidas para as Tabela 6.8a, 6.8b e 6.11a.
A partir do Índice de Facilidade obtido, pode -se concluir que mais de dois terços das questões objetivas da prova foram considerados, pelo menos, Difícil : das 26 questões válidas, 19 foram classificadas como Difícil e três, como Muito difícil.
Uma questão foi classificada como Fácil, e outras três consideradas Médio.
Nenhuma questão foi classificada como Muito fácil.
Já quanto aos índices de discriminação das questões objetivas do Componente de Conhecimento Específico da prova, tem -se como resultado a seguinte classificação : cinco das 26 questões foram consideradas boas, enquanto outras três tiveram Índice de Discriminação Muito bom.
Assim, para um pouco menos de um terço das questões – 8 em 26 – os Índices de Discriminação foram Bom ou Muito bom.
Dentre as demais, 13 delas foram classificadas como Médio e cinco como Fraco, sendo 18, por conseguinte, a quantidade de questões nos dois patamares mais baixos de discriminação.
Constata -se, assim, que a prova – no que se refere ao Componente de Conhecimento Específico – possuía média capacidade de discriminar entre aqueles que dominavam ou não o conteúdo.
A Tabela 6.12b apresenta a distribuição das questões levando -se em conta simultaneamente a classificação dos dois índices.
Dentre as questões que alcançaram os maiores índices de discriminação, classificadas com Índice de Discriminação Muito bom, uma delas, a de número 23, a de índice mais elevado ( 0,45 ), classificada como Fácil quanto ao Índice de Facilidade, com 66 \% dos estudantes marcando a opção correta.
A questão de número 28 foi a mais difícil dentre as 26 questões específicas válidas, com baixo Índice de Facilidade, apenas 9 \% de acertos.
Essa questão apresentou poder discriminatório médio, 0,23.
As questões 15, 21, 24, 31 e 33, pelo critério ponto-bisserial, foram consideradas inadequadas, e, por isso, foram eliminadas do cômputo da nota final.
A Tabela 6.12b apresenta a distribuição das questões levando -se em conta simultaneamente a classificação dos dois índices.
Dentre as questões que alcançaram os maiores índices de discriminação, classificadas com Índice de Discriminação Muito bom, uma delas, a de número 23, a de índice mais elevado ( 0,45 ), classificada como Fácil quanto ao Índice de Facilidade, com 66 \% dos estudantes marcando a opção correta.
A questão de número 28 foi a mais difícil dentre as 26 questões específicas válidas, com baixo Índice de Facilidade, apenas 9 \% de acertos.
Essa questão apresentou poder discriminatório médio, 0,23.
As questões 15, 21, 24, 31 e 33, pelo critério ponto-bisserial, foram consideradas inadequadas, e, por isso, foram eliminadas do cômputo da nota final.
A título de exemplo das análises do comportamento das questões objetivas, o Gráfico 6.5b analisa a questão 23 do Componente de Conhecimento Específico.
Esta questão foi considerada pelas respostas dos estudantes avaliados na prova como uma questão com nível de facilidade Fácil, 0,66, ou seja, 66 \% dos estudantes assinalaram acertadamente a opção E, correspondente ao gabarito.
Como já comentado, seu Índice de Discriminação foi igual a 0,45, classificado como Muito bom, o maior valor de discriminação.
Neste gráfico, cada uma das cinco curvas representa o percentual de respostas em determinada alternativa da questão 23, em função do número de acertos dos estudantes nessa parte da prova, antes de possíveis eliminações de questões pelo critério do ponto-bisserial.
A alternativa correta E, representada no gráfico pela curva em vermelho, foi escolhida em maiores proporções pelos alunos com desempenho melhor nessa parte da prova.
Já as alternativas incorretas, também denominadas distratores, foram selecionadas, principalmente, por aqueles com notas mais baixas.
Observa -se que a soma não é 100 \% por causa das questões não respondidas ou com mais de uma opção marcada.
Aqueles com nota zero, na sua quase totalidade deixaram esta questão em branco ou marcaram mais de uma alternativa, comportamento considerado inválido.
A proporção de alunos que selecionou a resposta correta E aumenta gradativamente, chegando a atingir 100 \% para 15 acertos, sofrendo uma pequena queda nos 19 acertos e voltando a 100 \% para 20 acertos ou mais, enquanto a proporção dos que escolheram alternativas incorretas decai, a partir de um ou dois acertos, como função do número de acertos nessa parte da prova.
Os gráficos relativos às demais questões do Conhecimento Específico constam do Anexo I. 6.3 ANÁLISE DAS QUESTÕES DISCURSIVAS Esta seção apresenta estatísticas selecionadas e histogramas das Questões Discursivas de Formação Geral ( 6.3.1 ) e Conhecimento Específico ( 6.3.2 ).
São também apresentadas e comparadas as médias de subpopulações caracterizadas por Grande Região. 6.3.1 Componente de Formação Geral As análises dos resultados de desempenho dos estudantes de Química ( Licenciatura/Bacharelado ) nas duas questões discursivas relativas à Formação Geral encontram -se na Tabela 6.13 e no Gráfico 6.6, desagregadas por Grande Região, segundo Habilitação e Modalidade de Ensino.
Na Tabela 6.13, observa -se que a nota Média nesse conjunto de questões foi acima da obtida nas objetivas pelos estudantes de Licenciatura.
Estes estudantes obtiveram, em Formação Geral, Média 50,0 nas questões objetivas e 55,7 nas questões discursivas.
A situação se repete para os alunos de Bacharelado : média menor nas objetivas ( 56,4 ) do que nas discursivas ( 56,7 ).
Considerando -se agora os alunos de cursos presenciais de Licenciatura, observa -se que a nota Média nesse conjunto de questões foi acima da obtida nas objetivas.
Estes estudantes obtiveram, em Formação Geral, Média 50,1 nas questões objetivas e 55,7 nas questões discursivas.
Pode -se notar também que o Desvio padrão nesse conjunto de questões foi um pouco maior do obtido nas objetivas, 20,3 nas questões objetivas e 23,6 nas questões discursivas.
A maior Média foi obtida na região Sul ( 59,1 ), e a menor, na região Norte ( 52,7 ).
A Mediana dos alunos de cursos presenciais de Licenciatura, neste componente, foi 61,5.
A Mediana na região Norte foi 58,5, na região Nordeste foi 60,3, na Sudeste foi 63,5, na Sul 64,0, e na Centro-Oeste, 61,5.
A nota Máxima ( 97,5 ) foi a mesma nas regiões Nordeste e Sudeste, sendo 93,5 na região Norte, 96,0, na região Sul e 97,0 na região Centro-Oeste.
A nota Mínima ( 0,0 ) foi a mesma em todas as regiões do Brasil.
Também para os alunos de cursos a Distância de Licenciatura a nota Média nesse conjunto de questões foi acima da obtida nas questões objetivas pelo mesmo grupo de estudantes.
Estes estudantes obtiveram, em Formação Geral, Média 48,8 nas questões objetivas e 55,4 nas questões discursivas.
Pode -se notar que o Desvio padrão nesse conjunto de questões foi também maior do que o obtido nas objetivas, 20,5 nas questões objetivas e 23,7 nas questões discursivas.
A maior Média foi obtida na região Centro-Oeste ( 60,5 ), e a menor, na região Norte ( 51,7 ).
A Mediana dos alunos de cursos a Distância de Licenciatura, neste componente, foi 61,5.
Nas regiões Norte ( 52,0 ), Nordeste ( 60,3 ) e Centro-Oeste ( 60,5 ), a Mediana foi menor, e na região Sudeste ( 62,5 ) a Mediana foi maior que o valor para o Brasil.
Não houve registro de cursos de Educação a Distância na região Sul.
A nota Máxima ( 94,5 ) foi obtida na região Sudeste, sendo 87,0 na região Norte, 92,0 na Nordeste e 60,5 na região Centro-Oeste.
A nota Mínima ( 0,0 ) foi obtida nas regiões Norte, Nordeste e Sudeste.
No Centro-Oeste, foi 60,5.
Já para os alunos de cursos de Bacharelado, a nota Média nas questões objetivas foi pouco menor do que as das discursivas.
Estes estudantes obtiveram nas questões objetivas Média 56,4 e nas questões discursivas, 56,7.
Pode -se notar que o Desvio padrão nesse conjunto de questões foi maior ao do obtido nas objetivas : 20,4 nas questões objetivas e 24,3 nas questões discursivas.
A maior Média foi obtida na região Sudeste ( 57,7 ), e a menor, na região Norte ( 53,7 ).
A Mediana dos alunos de cursos de Bacharelado foi 62,5.
Nas regiões Norte ( 61,8 ), Nordeste, Sul e Centro-Oeste ( 61,5 em cada ), a Mediana foi menor que a do Brasil.
Na região Sudeste ( 63,0 ) foi maior.
A nota Máxima ( 97,5 ) foi obtida na região Sudeste, sendo 96,0 nas regiões Nordeste, Sul e Centro-Oeste. e 92,0 na região Norte.
A nota Mínima ( 0,0 ) foi a mesma em todas as regiões do Brasil.
Como já comentado, não há registro de cursos a Distância de Bacharelado.
O Gráfico 6.6 representa a distribuição das notas nas questões discursivas do Componente de Formação Geral.
A moda desta distribuição ocorre no intervalo ( 60 ; 70 ] para as duas Habilitações, Licenciatura e Bacharelado.
Destaca -se também o intervalo [ 0 ; 10 ] com uma proporção próxima de 10 \% do total de notas, para ambas as Habilitações, sendo que no intervalo [ 0 ; 10 ] inclui -se, além da nota zero, a frequência de alunos que deixaram este tipo de questão em branco.
Na sequência, os resultados verificados para cada uma das questões discursivas de Formação Geral serão apresentados, estabelecendo -se relações com os conteúdos abordados em cada uma delas.
Os comentários da Banca de docentes corretores, a respeito do observado na correção das respostas dos estudantes, suas impressões e conclusões serão apresentados junto à análise de cada questão.
Cumpre esclarecer que, tendo em vista que as questões discursivas de Formação Geral são padronizadas, ou seja, constam de todas as provas, os comentários da Banca são os mesmos para todas as carreiras acadêmicas, sendo direcionados a todos os estudantes que participaram do Enade/2017.
A seguir, serão analisados os desempenhos da Área de Química ( Licenciatura/Bacharelado ) nas duas questões discursivas de Formação Geral do Enade/2017, comparando -se os resultados obtidos com comentários para cada questão. 6.3.1.1 Análise de Conteúdo da Questão Discursiva 1 do Componente de Formação Geral Os dados de Química ( Licenciatura/Bacharelado ), obtidos a partir das respostas à questão 1, encontram -se na Tabela 6.14 e no Gráfico 6.7, desagregados por Habilitação e Modalidade de Ensino.
Nessa questão – de melhor desempenho dentre as duas de Formação Geral – os alunos dos cursos de Licenciatura, tiveram Média 56,9.
A maior Média para a questão 1 entre os alunos de cursos de Licenciatura foi obtida na região Sul ( 59,6 ), e a menor, na região Norte ( 55,2 ).
Já os alunos, dos cursos presenciais de Licenciatura, obtiveram uma Média ligeiramente mais baixa, 56,8 e os de cursos a Distância, valor um pouco mais alto, 57,0.
A maior Média para a questão 1 entre os alunos de cursos presenciais de Licenciatura foi obtida na região Sul ( 59,6 ), e a menor, na região Norte ( 55,0 ).
Entre os alunos de cursos a Distância de Licenciatura, a maior Média foi obtida na região Nordeste ( 59,0 ), e a menor, na região Centro-Oeste ( 50,0 ).
Quanto à variabilidade das notas, o Desvio padrão das notas foi 27,2 para os alunos de educação presencial e 26,1 para os de Educação a Distância.
Para os cursos presenciais de Licenciatura, as medianas de três regiões ( Norte, Nordeste e Centro-Oeste ) e do Brasil como um todo foram iguais ( 60,0 ).
A Mediana do Sudeste e do Sul foi 65,0.
Já para os cursos a Distância, a mediana do Brasil como um todo foi 60,0, a mesma dos cursos presenciais.
Na região Norte da Educação a Distância, a Mediana foi 62,5, na região Nordeste foi 60,0, na Sudeste foi 65,0 e na Centro-Oeste, 50,0.
Não houve registro de participantes na região Sul da Educação a Distância da Habilitação Licenciatura.
As notas máximas e mínimas da questão discursiva 1 foram, para ambas as Modalidades de Ensino, as mesmas para todas as regiões do Brasil, respectivamente, 100,0 e 0,0, com exceção nos cursos a Distância da nota Mínima da região Centro-Oeste ( 50,0 ) e das notas máximas das regiões Norte ( 90,0 ) e Centro-Oeste ( 50,0 ).
Entre os alunos de Bacharelado, a performance nesta questão também foi um pouco melhor do que na questão discursiva 2.
Os alunos dos cursos de Bacharelado obtiveram Média 56,8.
A maior Média para a questão 1 entre os alunos de cursos de Bacharelado foi obtida na região Norte ( 57,8 ), e a menor, na região Sul ( 54,2 ).
Quanto à variabilidade das notas, o Desvio padrão das notas foi 27,7, sendo o maior Desvio padrão os das regiões Sudeste e Centro-Oeste ( 27,8 ) e o menor, o da Nordeste ( 27,0 ).
Não houve curso a Distância na habilitação Bacharelado.
Portanto, os registros totais de Educação Presencial coincidem.
Para os cursos de Bacharelado, a Mediana do Brasil e a das regiões Sul e Centro-Oeste foram 60,0.
A Mediana da região Norte foi 62,5, e as do Nordeste e Sudeste foram 65,0.
As notas máximas e mínimas da questão discursiva 1 foram as mesmas para todas as regiões do Brasil, respectivamente, 100,0 e 0,0.
O Gráfico 6.7 mostra a distribuição das notas na questão discursiva 1 do Componente de Formação Geral.
Observa -se que a maior frequência corresponde ao intervalo ( 70 ; 80 ] para ambas as Habilitações.
Também é importante a proporção dos alunos que deixaram a resposta a essa questão em branco, em torno de 10 \% dos participantes. 6.3.1.2 Comentários sobre a correção de Conteúdo das respostas à Questão Discursiva 1 O enunciado era claro e indicava um tema de extrema relevância, qual seja, a incidência de doença grave integrante do rol das DST ( Doenças Sexualmente Transmissíveis ), que encaminha para a análise do direito fundamental à saúde, que tem assento constitucional ( CRFB, art. 6º. caput : São direitos sociais a educação, a saúde, a alimentação, o trabalho, a moradia, o transporte, o lazer, a segurança, a previdência social, a proteção à maternidade e à infância, a assistência aos desamparados, na forma desta Constituição ).
A questão focalizava, por meio de três textos, um tema social extremamente relevante, qual seja, a relação entre homem e mulher nas suas interações íntimas, no caso representadas pela vida sexual, e o reflexo das relações de poder veiculadas pela sociedade, que apontam, no mais das vezes, o predomínio do homem.
Tal posição de superioridade é justificada pelas convenções arraigadas que acarretam a educação dos homens e mulheres por padrões patriarcais, que redundam em posições machistas.
Um dos textos demandava uma reflexão sobre a realidade dos casais que convivem em relações estáveis ou fugazes e não têm o controle dos efeitos das relações sexuais, que abrangem a gravidez não planejada e, em decorrência do não cuidado, a inserção de doenças transmitidas sexualmente.
Esse panorama é refletido na pouca influência das mulheres sobre os homens quanto à proteção, normalmente representada pelo uso de preservativos nas relações intimas, e após a constatação da doença, no encaminhamento das pessoas para tratamento.
Como efeito colateral, está a transmissão de doenças para os fetos.
Outro texto trazia à tona o maior cuidado que as mulheres têm, em geral, com sua saúde, notadamente no momento em que se descobre grávida e que redunda na realização de exames obrigatórios de pré-natal, que podem facilmente detectar doenças e, caso encontradas, ter o seu tratamento encaminhado.
A par disso, destaca -se que a crise de recursos financeiros agravou a situação, que já recebe influxos sociais negativos, como a dificuldade do homem em utilizar preservativos, realizar exames e procurar tratamentos quando adoece.
Como muito divulgado pelas mídias, a crise financeira do país repercutiu diretamente na prestação de serviços de saúde em todas as esferas : distrital, federal, estadual e municipal.
Assim, agravou -se o problema da falta de medicamentos essenciais para o tratamento da sífilis e, certamente, de outras doenças.
Como se verifica, a temática indicada no enunciado, representada pelos textos, pertence à realidade conhecida amplamente pelos estudantes, pois também é divulgada por todos os meios de comunicação.
Trata -se de tema geral que, no entanto, normalmente, não se inclui nas grades curriculares de cursos superiores que não sejam da área de Saúde.
Apesar disso, a temática é de conhecimento geral e relacionada aos direitos humanos, esses sim, ponto de integração com as diretrizes curriculares.
A linguagem utilizada era clara, indicando a existência dos fatos vinculados a questões relevantes e à saúde pública.
No entanto, como os textos trazem uma ampla gama de possibilidades de interpretação, ocorreram algumas compreensões diversas das esperadas.
Além disso, ressalta -se que a amplitude dos textos disponibilizados no enunciado fornece elementos que poderiam ser simplesmente reproduzidos pelos concluintes nas suas respostas.
O enunciado da questão encaminha a resposta para dois aspectos.
O estudante deveria abordar conteúdo pertinente à vulnerabilidade das mulheres às DST e o papel social do homem em relação à prevenção dessas doenças.
Para responder a esses aspectos, um número expressivo de concluintes utilizou somente os elementos dos textos motivadores, sem nada acrescentar.
A resposta deveria conter também uma segunda parte, a indicação de duas ações voltadas para o público masculino, a serem adotadas no âmbito das políticas públicas de saúde ou de educação, para a redução do problema.
Verificou -se que o desempenho dos estudantes nessa questão foi de médio para bom.
Ao se analisarem as respostas, constata -se que muitos foram os que apresentaram os tópicos presentes no padrão de resposta, amplo o suficiente para albergar as respostas apresentadas pelos concluintes.
Tal resultado pode ser justificado, tendo em vista tratar -se de tema bastante difundido nas diversas mídias, por ter sido apresentado de forma ampla pelos textos motivadores presentes no enunciado.
Por outro lado, o debate no ambiente universitário sobre as relações de poder entre homens e mulheres tende a ser marcante pelo movimento de empoderamento feminino, liderado por organizações sociais com ampla atuação nos estabelecimentos de nível superior.
De forma recorrente, as relações patriarcais entre homens e mulheres foram vinculadas ao modelo machista vigente.
Essas condições, segundo os concluintes, explicam a incidência da doença e a situação de vulnerabilidade das mulheres, bem como a função social dos homens.
Outras respostas apontaram o poder público, nas áreas de saúde e/ou educação, como órgãos responsáveis pela solução do problema.
Para os que se fixaram na área da saúde, foram frequentes a indicação da necessidade de mais medicamentos, vacinas, educação sexual, médicos e postos de atendimentos.
Alguns, no entanto, apenas indicaram ações individuais para o homem, e outros, ações individuais para as mulheres.
Também figuraram respostas que apresentavam ações da igreja como solução, além de outras que faziam referências a campanhas no seio das famílias. 6.3.1.3 Análise de Conteúdo da Questão Discursiva 2 do Componente de Formação Geral A Tabela 6.15 apresenta as estatísticas básicas das notas dos alunos na questão discursiva 2 da área de Química por Grande Região, segundo Habilitação e Modalidade de ensino.
Como já comentado, esta tabela mostra, em particular, que o desempenho médio dos estudantes de cursos de Química, tanto de Licenciatura, quanto de Bacharelado, na questão discursiva 2 foi inferior ao obtido na questão discursiva 1 : médias, respectivamente 56,9 e 56,8 na questão 1 e 53,9 e 56,2, na questão 2.
Entre os cursos de Licenciatura, a região Sul foi aquela cuja Média nessa questão foi maior ( 57,9 ), e a de menor Média foi a região Norte ( 48,9 ).
Quando se considera também a modalidade de ensino, a Tabela 6.15 mostra que o desempenho médio dos estudantes de cursos presenciais de Licenciatura na questão discursiva 2 ( Média 54,0 ) foi inferior ao obtido na questão discursiva 1 ( Média 56,8 ).
A região Sul foi aquela cuja Média nessa questão foi maior ( 57,9 ), e a de menor Média foi a região Norte ( 49,2 ).
Quanto à variabilidade das notas, o Desvio padrão de todo o Brasil foi 32,3.
O maior desvio nessa questão foi obtido na região Nordeste ( 33,5 ), enquanto o menor foi obtido na região Sul ( 30,4 ).
A Mediana para estes alunos de cursos presenciais de Licenciatura foi 70,0, a mesma das regiões Sudeste, Sul e Centro-Oeste.
Nas regiões Norte e Nordeste a Mediana foi 50,0.
As notas máximas ( 100,0 ) e as notas mínimas ( 0,0 ) foram as mesmas em todas as regiões, sem exceção.
A Média das notas dos estudantes de cursos a Distância de Licenciatura na questão discursiva 2 foi 52,7.
A região Centro-Oeste foi aquela cuja Média nessa questão foi maior ( 70,0 ), e a de menor Média foi a região Norte ( 42,8 ).
O Desvio padrão de todo o Brasil foi 31,9.
A Mediana para estes alunos de cursos a Distância de Licenciatura foi 50,0, o mesmo valor encontrado nas regiões Nordeste e Sudeste.
As notas máximas ( 100,0 ) e as notas mínimas ( 0,0 ) foram as mesmas em todas as regiões, com exceção das notas Mínima e Máxima da região Centro-Oeste que foram iguais a 70,0.
Não houve registro de cursos na região Sul de Educação a Distância na Habilitação Licenciatura.
Também não houve curso de Educação a Distância em nenhuma das Grandes Regiões na Habilitação Bacharelado.
A desagregação por Grande Região para as notas de Bacharelado na questão discursiva 2 mostra que a maior Média foi observada na região Sudeste ( 57,8 ) e a menor na região Norte ( 49,7 ).
O Desvio padrão das notas do conjunto destes alunos foi 32,9, sendo que o maior desvio nessa questão foi obtido na região Norte ( 37,1 ), enquanto o menor foi obtido na região Centro-Oeste ( 31,5 ).
A Mediana para estes alunos de cursos de Bacharelado foi 70,0, a mesma em três regiões : Norte, Sudeste e Sul.
As notas máximas ( 100,0 ) e as notas mínimas ( 0,0 ) foram as mesmas em todas as regiões.
O Gráfico 6.8 mostra a distribuição das notas na questão discursiva 2 do Componente de Formação Geral.
Observa -se que a maior frequência corresponde aos alunos que obtiveram nota dentro do intervalo ( 70 ; 80 ] para ambas as Habilitações.
Destacam -se, também, os alunos que deixaram a resposta em branco, com proporção próxima de 15 \% para ambas as habilitações. 6.3.1.4 Comentários sobre a correção de Conteúdo das respostas à Questão Discursiva 2 O enunciado era claro e abordava um tema de extrema relevância : a questão de gênero.
Trata -se de uma proposta polêmica, pois abrange a análise da inserção social de minoria vinculada à questão de gênero.
O enunciado solicitava que o concluinte discorresse sobre a importância do nome para as pessoas transgêneras e propusesse uma medida, no âmbito das políticas públicas, para facilitar o acesso dessas pessoas à cidadania.
A respeito da primeira parte, a importância do nome, esperava -se que o estudante abordasse as dificuldades que o uso do nome civil ( nome de registro, nome de batismo, nome da certidão de nascimento, por exemplo ) pode acarretar para as pessoas transgêneras, além de demonstrar conhecimento sobre a importância do nome social, direito já reconhecido através, principalmente, de decisões judiciais e atos administrativos, como decretos e portarias que dispõem sobre o uso do nome social.
Em atendimento à segunda parte do comando, o participante deveria propor uma política pública associada ao direito à cidadania das pessoas transgêneras.
Sobre esse aspecto, esperava -se a compreensão do significado de política pública, ou seja, propostas de âmbito governamental, e não, iniciativas de Organizações Sociais, igrejas ou grupos sociais de qualquer tipo.
O nascimento das pessoas, no Brasil, é objeto de inscrição em cartórios de Registro Civil das Pessoas Naturais, órgão responsável por estabelecer a certificação pública dos dados de todos os indivíduos que não somente nascem, mas casam, eventualmente se divorciam e vêm a óbito.
Daí o nome civil.
Para flexibilizar a correção, outras nomenclaturas foram consideradas, a saber : identificação de nascimento, nome de batismo, nome de registro, nome próprio, nome verdadeiro, que diferenciam o nome civil daquele que o indivíduo deseja utilizar para traduzir a sua opção pessoal de gênero, que se convencionou denominar nome social.
Quanto ao entendimento do termo transgênero, deve -se destacar que o movimento de mudança do gênero é contemporâneo e não está necessariamente vinculado à mudança de sexo.
Existem várias hipóteses em jogo.
Uma é a que envolve cirurgias de mudança de sexo ( a menina extirpa os seus seios e inclui órgão genital por intervenção cirúrgica ; o menino inclui seios e extirpa o seu órgão genital e inclui o órgão feminino também por intervenção cirúrgica ).
Há outro movimento que envolve apenas as vestimentas e o linguajar corporal, além dos relacionamentos ( as meninas assumem o trajar, o andar, os gostos dos meninos e mantêm relacionamentos afetivos com meninas ; os meninos assumem o trajar, o andar, os gostos das meninas e mantêm relacionamentos afetivos com meninos ).
Em relação à discussão da primeira solicitação do comando da questão, destaca -se que todos têm nome civil, que foi escolhido pelos seus pais ou responsáveis e, na falta deles, por autorização do Juiz de Direito, pautado no sexo detectado no nascimento.
Já o nome social é aquele que indica a preferência do indivíduo na sua identificação nas relações sociais e que destoa do seu gênero de registro civil.
Ressalte -se que o Brasil não possui ainda legislação específica sobre pessoas transgêneras.
Observa -se um relativo progresso no reconhecimento de direitos, através, principalmente, de decisões judiciais e atos administrativos, como decretos e portarias que dispõem sobre o uso do nome social.
Sem uma lei que defina os procedimentos da alteração dos documentos para pessoas transexuais, essa parcela da população LGBT ( lésbicas, gays, bissexuais, transexuais e transgêneros ) é obrigada a procurar na Justiça o reconhecimento de sua identidade, em processos que podem ser longos e que dependem do olhar de determinados, profissionais, como médicos, psicólogos, assistentes sociais, bem como do entendimento dos advogados, defensores públicos, juízes e do Judiciário, como um todo.
Atualmente, tramita na Comissão de Direitos Humanos da Câmara dos Deputados o Projeto de Lei João Nery ( 5002/2013 ), dos deputados Jean Willys ( PSOL-RJ ) e Erika Kokay ( PT-DF ), que determina que o reconhecimento da identidade de gênero seja um direito do cidadão.
O projeto recebeu o nome do primeiro trans-homem operado no Brasil.
Grande parte dos respondentes considerou a incongruência entre a imagem corporal, representada pelo gênero assumido pelos trans, e o nome que apresenta nos seus documentos como causa de sérios constrangimentos e embaraços e as dificuldades para que esses sujeitos acessem direitos básicos, como saúde, educação e empregabilidade em uma sociedade heteronormativa e preconceituosa.
Muitos demonstraram reconhecer que a identidade e expressão de gênero, através do nome social, pressupõe que sejam respeitadas a dignidade da pessoa humana, a intimidade, a liberdade e o direito à cidadania plena.
No entanto, também foram encontradas, em quantidades significativas, respostas que demonstravam a existência de uma interpretação social conservadora que não admite a mudança de sexo nem do nome social.
Existe um acentuado preconceito contra todos os movimentos que buscam identificação social diversa da civil.
Isso inclui os variados grupos : gays, travestis, transgêneros, dentre outros.
Já em relação ao segundo tópico a ser abordado pelo concluinte, é notório que as pessoas trans encontram -se destituídas de direitos no campo das políticas públicas, como Já em relação ao segundo tópico a ser abordado pelo concluinte, é notório que as pessoas trans encontram -se destituídas de direitos no campo das políticas públicas, como saúde, habitação, emprego, educação, justiça, dentre outros.
Nesse sentido, uma grande diversidade de medidas, no âmbito das políticas públicas, poderia ser proposta.
Ações voltadas para a capacitação de profissionais da educação, saúde, segurança pública, dentre outras áreas, para melhor atender às demandas e especificidades desse público alvo seriam necessárias.
Ainda em relação à educação, seriam também necessárias ações que contribuíssem para diminuir a evasão, pois o ambiente escolar é difícil para uma pessoa transgênera, por conta das situações discriminatórias e preconceituosas vindas dos colegas de classe, e mesmo das/dos professores e/ou de outros profissionais da área, não podem ser mais um complicador nesse processo.
Ainda no que tange ao assunto, percebe -se que, ao serem excluídas dos bancos escolares, as oportunidades no mercado de trabalho formal praticamente inexistem.
Um outro aspecto que se evidencia no cotidiano desse segmento da sociedade diz respeito às situações de violência sofridas pelas pessoas transgêneras.
Nesse sentido, algumas propostas foram elencadas para minimizar o problema e combater a violência.
Políticas envolvendo melhor formação de funcionários das polícias civil e militar, por exemplo, somariam esforços na qualificação das abordagens policiais e dos atendimentos prestados pelas forças de segurança, pois o respeito à orientação sexual e à identidade e expressão de gênero é, antes de tudo, o respeito à dignidade, à cidadania e à própria democracia.
Nesse sentido, revela -se essencial que, ao sofrerem agressões, sejam elas físicas, psicológicas ou morais no espaço público ou domiciliar, as pessoas trans possam ser atendidas, com respeito inclusive ao nome social, nas delegacias comuns e, inclusive, nas Delegacias Especiais de Atendimento à Mulher ( DEAM ), pois esse é o órgão institucional que possivelmente detém as qualificações necessárias para o atendimento ao gênero feminino ao qual as travestis e mulheres trans efetivamente pertencem.
Também ocorre uma latente inacessibilidade das pessoas trans à assistência jurídica.
Nesse sentido, é necessário implementar projetos de capacitação para profissionais do poder judiciário, como um todo, de modo a prestarem a assistência jurídica gratuita na defesa intransigente dos direitos sociais, civis, trabalhistas e políticos dessa população.
Dessa maneira, a Justiça agiria de forma mais acessível e sensível para o segmento, que terá suas especificidades observadas a começar pelo devido respeito do direito ao uso do nome social.
Nesse âmbito, o mais comum foi propor a adoção de medidas punitivas para quem viola o direito à autodeterminação de gênero e a elaboração de leis que garantam a mudança do nome e assegurem outros direitos para as pessoas trans.
Propostas envolvendo a formação de profissionais, como as mencionadas aqui, e de conscientização da população foram as mais frequentes.
Muitas outras enfatizavam o direito ao uso respeitado do nome social e a garantia de alteração do registro civil, independentemente da cirurgia de redesignação sexual ou transgenitalização.
Muitos questionaram o fato de não bastar a autodeclaração no processo de alteração de nome, enfatizando que o nome social é uma questão de princípio de dignidade e cidadania.
Apesar de pouco mencionada nas respostas corrigidas, políticas envolvendo a promoção da saúde são muito necessárias.
Cabe aos gestores de saúde ações direcionadas, não somente à prevenção e tratamento do HIV/Aids, mas também políticas de saúde que contemplem as especificidades dessa população, tais como o Ministério da Saúde, especialmente o SUS ( Sistema Único de Saúde ), incluir, de fato, o atendimento para pessoas trans para acompanhamento da terapia hormonal, que muitas demandam, respeitando -se o direito de utilização do nome social.
De modo geral, são muitas as políticas públicas a que precisam ser implantadas.
Urge a necessidade de uma agenda mais compromissada por parte dos gestores do nosso país com essa comunidade, que ainda é muito discriminada.
Logo, tais políticas devem contemplar suas necessidades básicas como o direito de acesso aos estudos, à capacitação para o mercado de trabalho, à profissionalização, ao acesso a bens e serviços de qualidade em saúde, habitação, segurança, cultura e assistência social e, certamente, o respeito ao uso do nome social, o que também contribui para essa inclusão.
Assim, as opções apresentadas no padrão de resposta funcionaram como exemplos de políticas públicas possíveis, já que seria difícil apresentar uma lista suficientemente ampla para abranger a grande variedade de respostas encontradas.
O desempenho dos estudantes nessa questão foi médio.
As respostas, de modo geral, indicaram a necessidade de uma atualização por parte dos concluintes quanto a temas essenciais e, no caso da presente questão, a análise sobre aspectos fundamentais da questão de gênero.
Muitos, possivelmente motivados pelos textos apresentados no enunciado, restringiram -se a discorrer sobre preconceito de gênero ou mesmo sobre o preconceito de modo mais geral, sem responder ao que foi solicitado no comando da questão.
Outros apresentaram discursos religiosos ou ideológicos, também sem tratarem do que foi solicitado. 6.3.1.5 Análise de Língua Portuguesa das Questões Discursivas do Componente de Formação Geral Os dados obtidos a partir das respostas às questões discursivas do Componente de Formação Geral de Química desagregados por Grande Região, segundo Habilitação e Modalidade de Ensino, no que tange à Língua Portuguesa, encontram -se na Tabela 6.16 e no Gráfico 6.9.
Nos cursos de Licenciatura em Química, a Média desse aspecto foi 57,2, ligeiramente abaixo do obtido pelos alunos de Bacharelado ( 57,9 ).
Quando se considera a modalidade de ensino, os alunos de cursos presenciais de Licenciatura obtiveram Média 57,2, o mesmo valor da Média total.
Com respeito à Língua Portuguesa, a Média foi um pouco maior para os alunos de cursos a Distância (57,5)..
A Mediana das notas de Língua Portuguesa para os cursos de Licenciatura foi 62,5 para a modalidade Presencial e 65,0 para a Distância.
Nos cursos de Bacharelado, a Mediana também foi 65,0.
A nota máxima não chegou a 100,0 em qualquer combinação de habilitação e modalidade.
O Gráfico 6.9 mostra a distribuição das notas de Língua Portuguesa do Componente de Formação Geral.
Observa -se que a maior frequência ( em torno de 30 \% ) corresponde aos alunos que obtiveram nota no intervalo ( 60 ; 70 ].
Destacam -se, também, os alunos que 6.3.1.6 Comentários sobre a correção das respostas de Formação Geral com respeito à Língua Portuguesa deixaram ambas as respostas em branco representando um pouco menos do que 10 \% do total, caracterizando -se como máximo local. 6.3.1.6 Comentários sobre a correção das respostas de Formação Geral com respeito à Língua Portuguesa Os enunciados apresentados em 2017 não explicitaram a exigência de elaboração de um texto “ dissertativo ”, como nos anos anteriores.
O participante deveria fazer uma exposição de seus conhecimentos sobre cada um dos assuntos e estruturar seus textos de acordo com as características do registro formal adequado à situação comunicativa – avaliação de conhecimentos.
Essa configuração determina exigências quanto aos seguintes aspectos : adequação da seleção vocabular, desenvolvimento do conteúdo, organização lógica das ideias, estruturação sintática dos períodos, utilização de procedimentos de encadeamento textual e de referenciação, obediência às exigências morfossintáticas próprias da modalidade escrita da norma-padrão, respeito às regras ortográficas e às regras de acentuação gráfica.
O padrão de resposta utilizado na avaliação das questões 1 e 2 considerou os aspectos relevantes ao bom desempenho linguístico como competências distintas, de modo a permitir um mapeamento detalhado do domínio dos recursos disponíveis na Língua Portuguesa para a comunicação escrita formal.
Com base nesse objetivo, foram avaliados os seguintes aspectos : ( a ) estruturação textual condizente com o gênero solicitado e o modo de organização textual expositivo adequado ao gênero – essa competência envolve : • estruturação sintática condizente com o padrão da modalidade escrita formal da língua portuguesa de modo a garantir a clareza necessária ; • distribuição do conteúdo do texto em parágrafos, de modo a garantir a sua organização temática ; • utilização de operadores discursivos que contribuam para a progressão temática do texto, estabelecendo relações lógicas entre as ideias apresentadas, tanto do ponto de vista intrafrasal, como do interfrasal ; • utilização de procedimentos de referenciação lexical e pronominal que permitam a retomada de referentes textuais ; • utilização de sinais de pontuação que contribuam para a organização lógica da frase e do texto ; • inteligibilidade relacionada ao atendimento das exigências de estruturação textual.
Espera -se, portanto, que o participante recorra a procedimentos linguístico-discursivos para organizar seu texto, permitindo o encadeamento lógico entre suas partes de forma a garantir a progressão e a coerência textuais.
Isso significa que os seguintes procedimentos foram considerados inadequados, de acordo com o padrão de resposta proposto : • elaboração de frases fragmentadas que comprometam a estrutura lógico-gramatical do texto ; • sequência justaposta de ideias sem encaixamentos sintáticos, reproduzindo hábitos da oralidade ; • elaboração de frase com apenas oração subordinada, sem oração principal ; • emprego equivocado do conector ( preposição, conjunção, pronome relativo, alguns advérbios e locuções adverbiais ) comprometendo a expressão da relação lógica entre duas ideias, com prejuízo da clareza do texto ; • emprego do pronome relativo sem a preposição, quando obrigatória ; • repetição ou substituição inadequada de palavras sem o emprego dos recursos oferecidos pela língua ( pronome, advérbio, artigo, sinônimo ) ; • emprego inadequado dos pronomes relativos “ cujo(a) ” e “ onde ” ;  utilização inadequada dos sinais de pontuação, comprometendo a clareza textual ;  inteligibilidade relacionada ao atendimento das exigências de estruturação textual. ( b ) respeito às convenções ortográficas da norma-padrão da Língua Portuguesa – essa competência envolve o domínio das regras de acentuação gráfica e da grafia padrão das palavras ( com ausência de abreviaturas próprias da linguagem da internet ), de acordo com as convenções estabelecidas pela legislação em vigor e consubstanciadas no Vocabulário Ortográfico da Língua Portuguesa, editado pela Academia Brasileira de Letras. ( Vale notar que, nessa edição de 2017, não foi aceita a legislação anterior, no caso das regras relativas ao uso do hífen e da acentuação gráfica ).
Espera -se que o participante :  grafe corretamente as palavras ;  respeite as regras de acentuação gráfica ;  empregue maiúsculas em início de frase, em nomes próprios de pessoas, lugares ou instituições ;  grafe as siglas com letras maiúsculas ou apenas com a primeira letra maiúscula quando formarem uma palavra ;  evite abreviações como p/, vc, tb, pq, tá, né, usadas muitas vezes em escrita informal e na internet ;  obedeça às regras de separação de sílabas no final da linha. ( c ) domínio dos diferentes aspectos morfossintáticos próprios da modalidade escrita formal da norma-padrão da Língua Portuguesa – essa competência envolve : a concordância nominal, a concordância verbal, a regência nominal, a regência verbal, a flexão nominal, a flexão verbal, a correlação entre os tempos verbais, a colocação pronominal e a utilização de sinais de pontuação que contribuam para a organização lógica da frase e do texto.
Espera -se que o participante :  flexione o verbo para estabelecer concordância de número com o sujeito da frase ;  flexione o artigo, o adjetivo e o pronome para concordar em número e em gênero com o substantivo a que se referem ;  observe a regência nominal e a verbal, utilizando a preposição adequada depois de um substantivo, um verbo ou um adjetivo ;  empregue adequadamente o acento grave indicador da crase entre uma preposição e um artigo ( a+a ) ;  obedeça às regras de colocação pronominal ( próclise e ênclise ), distintas dos hábitos da oralidade ou da escrita informal ( exigência de próclise com termo atrator, não exigência de que o pronome oblíquo se ligue ao verbo auxiliar por meio de hífen ) ;  flexione adequadamente verbos, substantivos, adjetivos e pronomes no que diz respeito à expressão das categorias gramaticais ;  flexione os verbos para expressar a correlação de modo e tempo nas estruturas subordinadas.
Com base nesses critérios, foram considerados como desvios de caráter morfossintático, e não como desvios ortográficos, as alterações que envolvem mudança de classe gramatical ou de forma flexional do verbo :  eliminação da marca de infinitivo ( -r- ) e substituição por acento agudo ou ausência total de marca do infinitivo ;  confusão entre “ ão ” e “ am ” nas formas verbais ;  confusão entre “ há ” e “ a ” ;  uso de hífen para separar pronome átono – tanto uso indevido quanto omissão ( exemplo : “ esperasse ”, em lugar de “ espera -se ” ; “ falar -mos ”, no lugar de “ falarmos ” ;  “ esta ” ( no lugar de “ está ” ) ; “ mais ” ( no lugar de “ mas ” ) ; “ e ” ( no lugar de “ é ” ).  verbos “ ter ” e “ vir ” que, na terceira pessoa do plural, não apresentarem o acento circunflexo, serão considerados como desvio de concordância, em aspectos morfossintáticos. ( d ) seleção vocabular adequada à modalidade escrita formal da Língua Portuguesa, exigida pela situação comunicativa – essa competência envolve a precisão na seleção/utilização do vocabulário relacionado à temática solicitada pela questão ; a ausência de marcas da oralidade, como termos de sentido muito genérico ( “ coisa ”, “ negócio ”, “ você ” ) e termos de registros mais informais ( como gírias, jargões, frases feitas, ditados populares, termos regionais ).
Assim, espera -se que o participante respeite a adequação vocabular não usando gírias ou expressões coloquiais, evite repetição desnecessária de palavras e utilize um vocabulário mais formal, como solicitado por um texto dissertativo.
Observações : - A inteligibilidade relacionada ao atendimento das exigências de estruturação textual foi avaliada na segunda competência ( relativa aos aspectos textuais ). - Os problemas de coerência textual provocados por uso indevido do vocabulário foram avaliados na terceira competência ( relativa aos aspectos morfossintáticos e vocabulares ). - Não foram considerados como desvios morfossintáticos os problemas de caligrafia ( -a/ -o, -s, -r ). - Não foram considerados os textos grafados integralmente em caixa alta, já que não seria possível distinguir alguma marcação especial para as letras em início de frase. - Cada desvio foi considerado como uma ocorrência, mesmo que dois desvios fossem relativos ao mesmo aspecto linguístico.
A grade de avaliação do desempenho linguístico considerou, portanto, três grandes grupos de competências, segundo os aspectos explicitados anteriormente : a ) domínio das convenções ortográficas : grafia de vogais e consoantes, uso de maiúsculas e minúsculas, emprego do hífen e acentuação gráfica ; b ) domínio dos procedimentos de estruturação textual do ponto de vista microestrutural : organização interna dos períodos, emprego de conectores para a articulação lógica entre os períodos e entre os parágrafos, emprego de marcas de referenciação lexical e pronominal ; utilização dos sinais de pontuação que contribuem para a organização lógica da frase ; c ) domínio das regras de caráter morfossintático estabelecidas como modelares do ponto de vista da modalidade escrita formal da norma-padrão da Língua Portuguesa : concordância nominal e verbal, regência nominal e verbal, colocação pronominal, flexão nominal e verbal, correlação entre tempos e modos verbais, ausência de marcas de oralidade.
A seleção vocabular adequada à modalidade escrita formal da Língua Portuguesa foi incorporada a essa última competência, tendo em vista a intersecção entre as duas do ponto de vista das exigências do registro formal da modalidade escrita da norma-padrão.
Apreciação geral do desempenho dos estudantes : A correção revelou desempenhos distintos dos participantes, marcados pela falta de repertório cultural da maioria dos participantes : - a questão 1 propiciou a oportunidade de maior desenvolvimento do tema solicitado, gerando consequências na estruturação textual, já que os textos foram mais longos, mais elaborados e fluentes.
Entretanto, em virtude da diversidade de abordagens dos textos motivadores ( a falta de penicilina, a recusa no uso da camisinha pelos homens, a maior frequência das mulheres na realização de exames ginecológicos, a transmissão das DST aos bebês durante a gestação ), observaram -se respostas que se configuraram como verdadeiras paráfrases, com comprometimento do caráter autoral desejável em uma questão discursiva de Formação Geral.
Alguns participantes selecionaram um trecho de cada texto motivador para compor suas respostas, outros copiaram trechos inteiros. - a questão 2 teve um comportamento distinto, devido ao caráter polêmico que o caracteriza.
Assim, os textos são mais curtos, com vocabulário repetitivo e preso aos textos motivadores.
Além disso, observou -se grande ocorrência de protesto em função do tema, ora por motivos religiosos ora por preconceito de gênero.
Quanto aos aspectos linguísticos analisados durante esta avaliação, os resultados observados estão descritos nos parágrafos seguintes.
Aspectos ortográficos : O desempenho dos participantes revelou uma diferença muito grande nos dois aspectos analisados nesta competência : baixo índice de desvios da grafia padrão e grande índice de desvios de acentuação.
Em vários casos, ocorre ausência completa de acentuação gráfica.
Os resultados revelam que a tendência dominante entre os universitários brasileiros é a eliminação da acentuação gráfica, talvez motivada pelos hábitos relacionados às redes sociais e pela ausência de esclarecimento dos meios de comunicação, das autoridades e das escolas sobre as decisões do Acordo Ortográfico da Língua Portuguesa de 1990.
Os casos mais sistemáticos de eliminação do acento indicador da sílaba tônica são :  palavras proparoxítonas ( “ sifilis ”, “ proximos ”, “ politicas ”, “ publicas ”, “ transgeneros ” ) ;  palavras paroxítonas terminadas em ditongo crescente ( “ necessario ”, “ noticiarios ”, “ individuo ”, “ dependencia ”, “ varios ”, “ propria ”, “ transmissiveis ”, “ ocorrencia ” ) ;  palavras paroxítonas com hiato ( “ saude ” )  palavras oxítonas ( “ ninguem ”, “ esta ”, “ ate ”, “ tambem ”, “ prevencao ” ).
Por outro lado, destaca -se o uso indevido do acento gráfico em palavras como “ genêro ”, “ melâncolia ”, “ prevênção ”, “ intervênção ”, por exemplo.
Há, também, uma tendência a acentuar paroxítonos como se fossem oxítonos : “ transmissivéis ”.
Quanto ao domínio das convenções relativas à grafia das palavras, observam -se desvios recorrentes, como a hipercorreção pela escolha de “ e ” no lugar de “ i ”, por influência de hábitos da oralidade ( “ descriminação ” no lugar de “ discriminação ”, “ entervenção ” no lugar de “ intervenção ” ) ; desvios de grafia relacionados à variação diastrática podem ser observados em “ estrupo ”, “ subjulgadas ”, “ precoseito ”, “ soubre ”, “ vecendo ”, “ indesencia ”, “ apolojia ”, “ fulga ”, “ dereitos ”, “ sifelis ”.
Observam -se, também, casos de inadequação no uso da maiúscula para destacar determinadas palavras-chave do texto, como “ Transgêneros ”, “ Brasileiros ”, “ Homens ”.
Destaque -se, também, o grande número de participantes que grafam os textos inteiramente em caixa alta.
Vale observar, também, que, ao contrário do que se esperava, não apareceram abreviaturas próprias do “ internetês ”, relacionadas ao uso de redes sociais e e-mails.
Aspectos textuais : Esta competência é a que se revela como a mais problemática entre os participantes, porque são muitos os problemas observados, desvios acumulados durante toda a formação do estudante e que não se resolvem com um estudo autodidata, como acontece com regras ortográficas ou morfossintáticas.
São eles : sequência justaposta de ideias sem encaixamentos sintáticos ; redução drástica de estruturas subordinadas, ao lado do aumento na frequência de estruturas coordenadas e absolutas ; redução no uso de conectores para expressar relações lógicas essenciais à construção do texto, substituídas pela exigência de inferência por parte do interlocutor para suprir a sua ausência ; emprego equivocado de operadores que não estabelecem relações lógicas coerentes entre ideias do texto ; emprego inadequado do pronome relativo ( com omissão da preposição ou a utilização de pronome inadequado, como “ onde ” ) ; repetição exaustiva de termos sem a utilização de procedimentos mais sofisticados de substituição ( hiperonímias, hiponímias, nominalizações, expressões metafóricas ) ; frases fragmentadas que comprometem a estrutura lógico-gramatical ; frases formadas apenas por oração subordinada, sem oração principal.
Um importante aspecto a destacar é o baixo desempenho de uma parte dos participantes em relação à estrutura formal do texto produzido, o que é preocupante ao se levar em conta que são graduandos em fase final de formação.
São frequentes os casos de desvios de estruturação frasal, com uso inadequado ou ausência de conectivos entre parágrafos e entre frases.
Em uma parte dos textos, falta textualidade e domínio do registro padrão da língua.
Na verdade, observam -se relações linguísticas quase agramaticais, como Aspectos morfossintáticos e vocabulares : Em relação à regência, o desvio mais frequente é a falta do sinal indicativo da crase – isso revela que o usuário não tem consciência de que, sob a forma do termo “ a ”, existe a presença de uma contração entre a preposição “ a ” ( exigida pela regência do termo anterior ) e o artigo definido “ a ”.
Um desvio de regência significativo, nos últimos anos, é a utilização as estabelecidas pela sequência de gerúndios sem o apoio de um ponto de partida para a organização das informações gramaticais e semânticas.
Observou -se que uma parte dos participantes não distribuiu as ideias em parágrafos, talvez devido ao pequeno número de linhas disponibilizadas para a resposta da questão ou, quem sabe, pela suposição de que não seria necessária essa divisão por não se tratar de um texto no modelo de uma redação dissertativo-argumentativa, como solicitado nos vestibulares.
Em função do encaminhamento dos enunciados das duas questões, que solicitaram uma análise do problema e encaminhamentos de políticas públicas, houve uma grande tendência, também, de construção de dois parágrafos desconexos, sem utilização de elementos coesivos adequados à progressão textual.
Quanto à utilização dos mecanismos de referenciação, deve -se destacar a ocorrência de repetições de palavras ou expressões sem a utilização de termos sinônimos ou pronomes, como seria adequado.
Quanto à utilização dos sinais de pontuação, observou -se uma grande precariedade nos textos analisados.
É muito frequente a ocorrência de parágrafos sem marca interna de pontuação para separar os períodos.
Vale observar que não foi penalizada a ausência de vírgula para destacar locuções ou adjuntos adverbiais de pequena extensão deslocados de posição na frase, por ser um uso opcional.
São os seguintes os tipos de problemas encontrados : a ) vírgula : utilização de vírgula para separar o sujeito e o predicado ; ocorrência de apenas uma das vírgulas para separar uma palavra, uma expressão ou uma oração encaixada ; uso de vírgula no lugar do ponto para separar ideias que constituem períodos distintos ; ausência de vírgula para separar elementos de uma enumeração ; ausência de vírgula para separar oração adjetiva explicativa ou utilização inadequada para separar oração adjetiva restritiva ; b ) ponto e vírgula : utilização do ponto e vírgula no lugar de vírgula ; c ) ponto final : ausência de ponto final para separar períodos.
Aspectos morfossintáticos e vocabulares : Em relação à regência, o desvio mais frequente é a falta do sinal indicativo da crase – isso revela que o usuário não tem consciência de que, sob a forma do termo “ a ”, existe a presença de uma contração entre a preposição “ a ” ( exigida pela regência do termo anterior ) e o artigo definido “ a ”.
Um desvio de regência significativo, nos últimos anos, é a utilização inadequada de uma preposição ou sua ausência após o verbo ou o nome ( substantivo ou adjetivo ).
Outro problema relacionado à regência verbal e à nominal, encontrado frequentemente nas questões, foi a ausência de preposição antes de pronome relativo, processo generalizado na modalidade oral da língua, em situações de registro informal.
Apesar da possibilidade de que essa alteração de regência se generalize no padrão escrito da Língua Portuguesa, como já está ocorrendo até em textos jornalísticos, o não emprego da preposição foi considerado inadequado neste processo de avaliação.
A concordância verbal e a concordância nominal apresentam alguns desvios muito frequentes.
Quanto à concordância de número, observou -se ausência de marca ( com sujeito anteposto ou posposto ) ou uso indevido ( uso inadequado da marca de plural comandado pelo núcleo plural da locução adjetiva, apesar de o substantivo que funciona como núcleo do sintagma nominal estar no singular ).
Uma ocorrência que se destacou foi a ausência de acento circunflexo na forma plural do presente do indicativo do verbo “ ter ”, que foi considerada como um desvio na concordância verbal e não na acentuação gráfica.
Quanto à concordância de gênero, vários casos foram observados, normalmente no âmbito de sintagmas nominais longos, em que o adjetivo está afastado do substantivo.
Deve -se destacar o aparecimento da marca de plural em verbos ou adjetivos comandados por adjunto adnominal plural, apesar de serem relacionados a núcleos substantivos no singular, evidenciando um processo de hipercorreção ( exemplo : “ O nome social dos trangêneros podem ser estabelecidos por uma legislação específica ” ).
Quanto à questão da colocação pronominal, foram poucos os desvios observados.
Concluiu -se que, no registro escrito formal, a maioria dos participantes já incorporou regras como a não introdução da frase por um pronome oblíquo e a próclise na presença de um termo atrator.
Não se adotou, entretanto, o padrão excessivamente formal descrito pelas gramáticas normativas em relação à posição do pronome oblíquo em locuções verbais, já que esse uso está muito distante da prática cotidiana, até em textos mais formais.
Quanto aos aspectos vocabulares, alguns tipos de inadequação foram observados : expressões da oralidade ; seleção vocabular incompatível com o contexto, gerando falta de inteligibilidade ; falta de domínio de vocabulário mais abstrato e de maior complexidade, essencial ao desenvolvimento do texto de base dissertativa.
O principal aspecto observado foi a excessiva repetição de certas palavras, revelando limitação de repertório vocabular.
O termo “ pessoa ”, por exemplo, chega a ser repetido até 7 ou 8 vezes em um mesmo texto.
Várias marcas de oralidade foram identificadas, embora não com alta frequência : uso do pronome relativo “ onde ” como relativo universal, falta de artigo definido antes de substantivo, reduções como “ tá ”, “ pra ”, “ pro ”, “ prum ”, expressões informais, eliminação de preposições. substantivo, reduções como “ tá ”, “ pra ”, “ pro ”, “ prum ”, expressões informais, eliminação de preposições.
Em função do tema solicitado na questão 1, a sigla DST foi grafada de diferentes maneiras : DST´s, DSTs ou DSTS.
Além disso, houve flutuação de gênero gramatical nos sintagmas que envolviam a sigla : “ os diferentes DSTs ” ou “ as diferentes DSTs ”.
Em função do tema solicitado na questão 2, houve flutuação de gênero gramatical nos sintagmas que envolviam a sua palavra-chave : “ os trans ” ou “ as trans ” ; “ os transgêneros ” ou “ as transgêneras ”.
Todas essas formas foram aceitas, por considerarmos que sua utilização é muito recente e os participantes ainda não sistematizaram essas formas em seu uso da língua. 6.3.2 Componente de Conhecimento Específico As estatísticas básicas das notas referentes às questões discursivas do Componente de Conhecimento Específico de Química Licenciatura por Grande Região segundo Modalidade de Ensino estão disponibilizadas na Tabela 6.17a.
As informações equivalentes para o Bacharelado estão na Tabela 6.17b.
Observa -se que as Médias foram bem mais baixas do que para as questões discursivas do Componente de Formação Geral.
Enquanto no Componente de Formação Geral as Médias para estudantes de Química de Licenciatura foram, respectivamente para os de Ensino Presencial e a Distância, 55,7 e 55,4 ( 55,7 para o total ), na parte de Conhecimento Específico, as Médias foram 10,9 e 8,0 ( 10,7 para o total ).
Já para os cursos de Bacharelado, a Média foi 56,7 para as discursivas de Formação Geral ( não houve Educação a Distância no Bacharelado, portanto, esta nota é a Média tanto da Educação Presencial quanto total ) e 22,2 para as de Conhecimento Específico.
Para os cursos de Licenciatura, a nota Máxima foi maior para os cursos presenciais ( 91,7 ) do que para os a Distância ( 73,3 ).
Para os cursos de Bacharelado a nota Máxima foi 91,7 e a nota Mínima foi zero.
O Gráfico 6.10 representa a distribuição das notas nas questões discursivas no Componente de Conhecimento Específico.
A moda desta distribuição ocorre no intervalo [ 0 ; 10 ], com frequência de cerca de 30 \% dos participantes da Habilitação Bacharelado e mais do que 60 \% da Habilitação Licenciatura.
Observa -se que nesse intervalo estão computadas as ocorrências de respostas em branco para as três questões discursivas de Conhecimentos Específicos.
A análise de cada uma destas questões será feita a seguir. 6.3.2.1 Análise da Questão Discursiva 3 do Componente de Conhecimento Específico A questão 3 foi igual para os estudantes das duas habilitações.
Nessa questão, cujos resultados aferidos encontram -se descritos nas Tabelas 6.18a e 6.18b, respectivamente para os cursos de Licenciatura e de Bacharelado, as Médias dos estudantes de cursos de Licenciatura foram 9,3 e 7,0, respectivamente para os cursos presenciais e os a Distância ( 9,1 para o total ), ao passo que a Média dos estudantes de Bacharelado foi 27,0 ( a Média da Educação Presencial coincide com a total, haja vista que não há curso de Educação a Distância ).
Esta foi a questão de melhor desempenho dentre as três discursivas de Conhecimento Específico para o Bacharelado e a de desempenho intermediário para a Licenciatura.
Para ambas as habilitações, a nota Máxima foi 100,0 e a Mínima foi 0,0 em todas as regiões do Brasil, com exceção na região Norte na Licenciatura.
O Gráfico 6.11 mostra a distribuição das notas na questão discursiva 3, do Componente de Conhecimento Específico da área de Química ( Licenciatura/Bacharelado ).
Desconsiderando que, na Licenciatura, quase 30 \% dos estudantes deixaram a resposta à questão em branco e mais de 40 \% receberam nota zero, e que, no Bacharelado, quase 20 \% deixaram a resposta à questão em branco e mais de 20 \% receberam nota zero, a moda dessa distribuição seria o intervalo ( 10 ; 20 ], para ambas as habilitações. 6.3.2.2 Comentários sobre as respostas à Questão Discursiva 3 A Questão 3 foi apresentada para estudantes de Bacharelado e Licenciatura em Química.
Abordou conteúdos fundamentais do curso de Química, referentes a variáveis termodinâmicas e espontaneidade de reações químicas, conteúdos do campo de Físico-Química.
Esses conteúdos são fundamentais na formação em Química e estão alinhados às Diretrizes que norteiam os cursos de graduação.
O nível de profundidade esperado pelo padrão de respostas pode ser considerado médio.
O enunciado consistiu em um texto bem elaborado, claro e com linguagem adequada, não havendo possibilidade de interpretações diferentes da esperada.
Em comparação com as demais questões de formação específica, a Questão 3 envolvia conhecimentos simples e básicos, acarretando um grau de profundidade menor do que nas questões discursivas 4 e 5.
Como mencionado, a questão envolveu a abordagem de três variáveis termodinâmicas, variação de entalpia ( H ), variação de entropia ( S ) e variação da energia livre de Gibbs ( G ), e sua correlação em uma equação termodinâmica que possibilita a predição da espontaneidade de uma reação química.
O padrão de respostas foi amplo o suficiente para contemplar interpretações parciais acerca de uma ou duas dessas variáveis sem a completa conclusão da resposta.
Analisando -se as respostas apresentadas, constatou -se um baixo índice de respostas inteiramente corretas.
Na maioria dos casos, observou -se que os estudantes conseguiram analisar ao menos um dos parâmetros termodinâmicos ( G, H e S ).
As principais interpretações corretas foram : - H \&lt ; 0 por ser uma reação exotérmica. - G \&lt ; 0 por ser um processo espontâneo. - S \&lt ; 0, pois Sgás \&gt ; Sliq, e o reagente é formado apenas por gases, enquanto no produto há predomínio de líquidos.
Os estudantes que responderam corretamente à questão, ainda que de forma incompleta, mostraram domínio da linguagem científica e compreensão dessas variáveis.
Não houve uma tendência à incorreção completa das respostas, mas os erros foram os mais diversos.
Um tipo de resposta incorreta encontrada com certa frequência foi aquela na qual o concluinte apenas parafraseou o enunciado, sem acrescentar contribuições e desenvolver um raciocínio lógico.
Outro aspecto a ser considerado foi que, mesmo conseguindo identificar os sinais de cada variável termodinâmica, vários concluintes demonstraram -se incapazes de estabelecer um raciocínio matemático envolvendo as variáveis para concluir que, para se ter G \&lt ; 0 ( condição de espontaneidade ), tem -se que lHl \&gt ; lTSI.
As principais fontes de erros que levaram a esses baixos resultados foram o desconhecimento do tema e das variáveis termodinâmicas envolvidas e a dificuldade em se construir um raciocínio matemático, a partir da análise de uma equação termodinâmica apresentada no enunciado.
Um baixo percentual de concluintes alcançou um resultado satisfatório.
Da análise das respostas, é possível concluir que, tanto os concluintes de Bacharelado quanto os concluintes de Licenciatura, tiveram, em média, um baixo desempenho.
Os concluintes de Bacharelado demonstraram um melhor conhecimento do tema que se refletiu em respostas mais consistentes.
Esse fato foi evidenciado pelas médias obtidas.
Enquanto a média dos graduandos de Bacharelado foi 33,5, a dos de Licenciatura foi de 13,0.
A subdivisão da pontuação no sentido de considerar a abordagem de apenas uma ou de duas variáveis termodinâmicas possibilitou um escalonamento das notas e a diferenciação entre os níveis de respostas apresentadas.
De forma geral, as respostas que receberam notas mais baixas envolveram no máximo a compreensão correta de uma das variáveis termodinâmicas envolvidas na questão.
As respostas medianas, em geral, englobaram a compreensão de duas das variáveis termodinâmicas.
As respostas boas foram aquelas que abordaram corretamente as três variáveis termodinâmicas e demonstraram um raciocínio matemático correto.
Ao se analisar uma questão que abordou um conteúdo básico na formação na área de Química e, ainda com um enunciado que incluiu equação termodinâmica ( fato esse que descarta erros decorrentes de memorização ), constatou -se o completo desconhecimento do tema ou a dificuldade da expressão do conhecimento na construção de um texto discursivo. 6.3.2.3 Análise da Questão Discursiva 4 do Componente de Conhecimento Específico As Tabelas 6.19a e 6.19b contêm as informações relativas à questão 4 do conjunto de questões do Componente de Conhecimento Específico.
A questão 4 também foi a mesma para as duas habilitações.
Nos cursos de Licenciatura, o desempenho dos estudantes nessa questão foi ligeiramente superior ao da questão 3 e inferior ao da questão 5.
Nos cursos de Bacharelado, o desempenho foi o menor dentre as três questões.
As Médias para os alunos de Licenciatura ( ver Tabela 6.19a ) foram, respectivamente para os de cursos presenciais e a Distância, 9,2 e 9,3 ( 9,2 para o total ).
A nota Máxima ( 100,0 ) foi atingida por, pelo menos, um concluinte das regiões Sudeste e Centro-Oeste na modalidade presencial e por nenhum estudante da modalidade a Distância.
A Mediana foi zero para ambas as modalidades de ensino e para todas as regiões, com exceção da região Sudeste na modalidade presencial e da região Centro-Oeste na modalidade a Distância, que obtiveram Mediana 10,0.
A Mediana zero indica que, pelo menos, metade dos participantes obtiveram nota zero.
A nota Mínima foi zero para todas as combinações de modalidade de ensino e região, com exceção da região Centro-Oeste na modalidade a Distância, que obteve nota Mínima igual a 10,0.
A Média para os alunos de Bacharelado ( ver Tabela 6.19b ) foi 16,4.
A nota Máxima ( 90,0 ) foi atingida por, pelo menos, um concluinte das regiões Nordeste e Sudeste.
Nas regiões Sul e Centro-Oeste, a nota Máxima foi 70,0, e na Norte, 50,0.
A Mediana foi 10,0 para o Brasil e para quatro regiões, a exceção foi a Norte, cuja Mediana foi 5,0.
A nota Mínima foi zero para todas as regiões.
O Gráfico 6.12 representa a distribuição de notas da questão discursiva 4, do Componente de Conhecimento Específico.
Essa distribuição tem moda na classe de estudantes que, tendo respondido à questão, obtiveram nota zero para ambas as Habilitações, com frequência em torno de 22 \% para o Bacharelado e de cerca de 32 \% para a Licenciatura. 6.3.2.4 Comentários sobre as respostas à Questão Discursiva 4 A Questão 4, comum aos estudantes de Bacharela e de Licenciatura, apresentou um enunciado bem elaborado, claro e com linguagem adequada.
A questão estava baseada em um contexto histórico da Ciência e foi dividida em três itens, o que favoreceu a abordagem de diferentes conteúdos de Físico-Química : termoquímica, cinética química e equilíbrio químico.
Os conteúdos abordados são fundamentais na formação em Química e estão alinhados às Diretrizes que norteiam os cursos de graduação.
O nível de profundidade exigido pode ser considerado de médio para alto e não havia possibilidade de interpretações diferentes da esperada.
Em comparação com as demais questões de conhecimentos específicos, esta questão envolvia conhecimentos fundamentais de diferentes conteúdos, o que aumentou seu grau de dificuldade, especialmente em relação à Questão 3 que também avaliou conhecimentos de Físico-Química.
O padrão de respostas apresentou -se claro e abrangente, possibilitando a pontuação para respostas com desenvolvimento parcial.
Como dito anteriormente, a questão foi subdivida em três itens.
O primeiro item, ‘ a ’, solicitava a explicação da importância da obtenção de amônia a partir do nitrogênio atmosférico.
Em geral, as respostas associavam a necessidade de se produzir amônia para seu uso na produção de insumos bélicos e agrícolas.
Entretanto, foram pouquíssimas as respostas que indicaram a amônia como matéria-prima para a produção de nitrato, o qual era usado na produção bélica e agrícola, em substituição ao salitre.
A grande maioria apontava a amônia como um simples substituto do salitre.
Outras justificativas utilizadas envolveram o uso do nitrogênio atmosférico em substituição à extração mineral do salitre, devido à grande disponibilidade do nitrogênio atmosférico, que o faz uma matéria-prima abundante na rota industrial economicamente viável para a síntese da amônia.
No item ‘ b ’, o raciocínio comum foi ao encontro do proposto no padrão de respostas.
Por ser uma reação exotérmica, o aumento da temperatura favorece o sentido inverso, que é endotérmico.
Assim, o equilíbrio é deslocado no sentido dos reagentes, e a concentração de amônia diminui.
Esse item foi o que apresentou as melhores respostas, sem fuga ao proposto pelo padrão.
Cabe lembrar que a variação de entalpia também foi abordada na Questão 3.
O item ‘ c ’ demonstrou ser o de maior dificuldade para os estudantes.
A maioria das respostas apresentadas indicava que o aumento da temperatura acarreta o aumento da velocidade da reação.
Entretanto, poucos concluíram corretamente o raciocínio, indicando que o aumento da velocidade em ambos os sentidos faz com que se alcance mais rapidamente a condição de equilíbrio da reação e que a contínua remoção de amônia desloca o equilíbrio no sentido de aumentar sua produção.
Algumas justificativas equivocadas não foram consideradas, tais como : a maior temperatura volatiliza os componentes ; é mais economicamente viável trabalhar -se em maiores temperaturas ; a maior temperatura degrada a amônia ; pode trabalhar -se em pressões menores em altas temperaturas ; ou que é mais fácil se trabalhar em temperaturas mais altas industrialmente.
Tal qual observado na Questão 3, analisando -se as respostas corretas, verificou -se que os concluintes mostraram domínio da linguagem científica e compreensão dessas variáveis.
No entanto, para os três itens, verificaram -se respostas insuficientes nas quais o concluinte apenas parafraseava o enunciado, sem acrescentar contribuições e desenvolver um raciocínio lógico.
Um baixo percentual de concluintes alcançou um resultado satisfatório.
Da análise das respostas, foi possível concluir que tanto os graduandos de Bacharelado quanto os graduandos de Licenciatura tiveram, em média, um desempenho muito aquém do esperado.
Os graduandos de Bacharelado obtiveram média igual a 19,9, e os de Licenciatura alcançaram média igual a 13,1.
As principais fontes de equívocos que levaram a esse resultado foram : - No item ‘ a ’, o desconhecimento de que a amônia não é um substituto do salitre, mas se trata da matéria-prima para a produção de nitratos, estes sim os substitutos do salitre.
Cabe destacar que o termo salitre do Chile, apesar de popular de grande importância histórica, não corresponde a uma nomenclatura oficial.
Este fato pode ter sido o responsável pela incompreensão de muitos concluintes. - No item ‘ c ’, a principal dificuldade foi a correlação entre os conceitos de velocidade de reação e equilíbrio químico.
De forma geral, as notas mais baixas foram de respostas com acerto de apenas um dos itens, geralmente o item ‘ a ’ ou o item ‘ b ’.
As respostas medianas apontaram a compreensão de dois dos itens, novamente com destaque para os itens ‘ a ’ e ‘ b ’.
Já as respostas boas foram aquelas nas quais os estudantes conseguiram desenvolver respostas satisfatórias nos três itens, seja construindo um texto correto para os três itens ou ainda correto em dois itens e parcialmente correto em um terceiro.
A baixa média alcançada reflete o reduzido percentual de concluintes que conseguiram desenvolver uma resposta satisfatória.
Esse resultado é preocupante na medida em que a questão abordou conteúdos básicos na formação na área de Química, envolvendo conhecimentos de três temas importantes : termoquímica, cinética química e equilíbrio químico, bem como um contexto histórico.
Além de possíveis lacunas na formação inicial dos graduandos, também foi possível observar a dificuldade na expressão do conhecimento adquirido por meio de um texto discursivo.
Comparando -se os resultados das questões 3 e 4, observou -se que as médias dos graduandos de Licenciatura foram praticamente iguais ( 13,0 e 13,1 ), enquanto houve uma queda na média dos graduandos de Bacharelado de 33,5 para 19,9. 6.3.2.5 Análise da Questão Discursiva 5 do Componente de Conhecimento Específico A questão discursiva de número 5, diferentemente das duas anteriores, não foi igual para as duas habilitações.
Para o Bacharelado, a questão abordava procedimentos e normas de segurança para situações que podem ocorrer no dia a dia de um laboratório.
Para os estudantes de Licenciatura, a questão 5 abordava uma metodologia de ensino de Química, o uso de analogias.
As Tabelas 6.20a e 6.20b contêm, respectivamente para Licenciatura e Bacharelado, as informações relativas à questão 5 do Componente de Conhecimento Específico.
O desempenho dos estudantes da Licenciatura de ambas as habilitações nestas questões foi o mais alto dentre as três questões discursivas desse componente.
Para os estudantes de Bacharelado, o desempenho foi mais alto do que o da questão 4 e mais baixo do que o desempenho na 3.
As Médias para os alunos de Licenciatura ( ver Tabela 6.20a ) foram, respectivamente para os de cursos Presenciais e a Distância, 14,3 e 7,8 ( 13,8 para o total ).
A nota Máxima ( 100,0 ) foi atingida por, pelo menos, um concluinte de todas as regiões na modalidade Ensino Presencial.
Os estudantes do Ensino a Distância, Licenciatura, obtiveram nota Máxima 75,0.
A Mediana foi 0,0 nas duas modalidades de ensino da Licenciatura e em todas as regiões.
Cumpre lembrar que para aquelas combinações de modalidade e região com Mediana zero, temos que pelo menos metade dos alunos obtiveram esta nota.
A nota Mínima foi zero para todas as combinações de modalidade de ensino e região.
A Média para os alunos de Bacharelado ( ver Tabela 6.20b ) foi 23,3.
A nota Máxima ( 90,0 ) foi atingida por, pelo menos, um concluinte das regiões Sudeste e Sul.
Na região Norte, a nota Máxima foi 75,0, na Nordeste foi 65,0 e na região Centro-Oeste foi 75,0.
A Mediana do Brasil foi 20,0, a mesma da região Sudeste.
No Norte, Nordeste e Centro-Oeste, a Mediana foi 15,0 e no Sul foi 25,0.
A nota Mínima foi zero para todas as regiões.
Cumpre salientar, mais uma vez, que a habilitação Bacharelado não apresentou participante de modalidade Ensino a Distância, o que faz com que os registros feitos para a habilitação Ensino Presencial coincidam com os registros totais.
O Gráfico 6.13 representa a distribuição das notas da questão discursiva 5 do Componente de Conhecimento Específico.
Essa distribuição tem moda nas notas zero para a Licenciatura e no intervalo ( 20, 30 ] para o Bacharelado.
Este intervalo, ( 20, 30 ], também é o segundo de maior frequência para a Licenciatura.
O Gráfico 6.13 representa a distribuição das notas da questão discursiva 5 do Componente de Conhecimento Específico.
Essa distribuição tem moda nas notas zero para a Licenciatura e no intervalo ( 20, 30 ] para o Bacharelado.
Este intervalo, ( 20, 30 ], também é o segundo de maior frequência para a Licenciatura. 6.3.2.6 Comentários sobre as respostas à Questão Discursiva 5 Questão 5 – Bacharelado Essa questão foi específica para os formandos de Bacharelado e abordou procedimentos e normas de segurança em três situações que podem ocorrer no dia a dia de um laboratório.
Cada uma dessas situações correspondeu a um item específico no enunciado da questão.
Os conteúdos abordados são fundamentais na formação em Química e estão alinhados às Diretrizes que norteiam os cursos de graduação.
O nível de profundidade exigido pode ser considerado variando de médio a alto e não havia possibilidade de interpretações diferentes da esperada.
A questão apresentou um enunciado bem elaborado, claro e com linguagem adequada.
O padrão de respostas apresentou -se claro e abrangente, possibilitando a pontuação para respostas com desenvolvimento parcial.
Em comparação com as demais questões de formação específica, esta questão envolvia conhecimentos práticos de procedimentos de segurança em laboratório.
Mas esses conhecimentos não são empíricos e derivam da aplicação de conceitos fundamentais de Química em situações específicas.
Como já comentado, a questão estava subdivida em três itens.
O item ‘ a ’ foi aquele com maior percentual de acertos, mesmo que apenas parciais.
A resposta mais comum foi interpretar a dissolução do ácido em água como um processo exotérmico, sem um detalhamento do papel da água como arrefecedor.
A projeção de ácido, como consequência da não adoção do procedimento, também foi citada como resposta frequente.
No item ‘ b ’, a resposta mais comum foi a de acúmulo do gás no interior da mangueira.
Entretanto, a maioria das consequências apontadas remeteu a incêndios e a explosões do gás, como se houvesse um vazamento de gás no laboratório e não apenas a liberação da quantidade presente na mangueira.
Esse tipo de resposta só foi aceito quando o concluinte justificou o vazamento de maior volume de gás por possível falha na válvula de alimentação.
O item ‘ c ’ foi o de menor percentual de acertos.
O refluxo da água foi uma explicação comum, geralmente justificada por uma diferença de pressão entre o interior do kitasato e a mangueira.
Como consequência equivocadamente citada, a principal resposta foi a quebra da vidraria pela diferença de pressão.
Como as três operações citadas no enunciado remetem à segurança de laboratórios, muitas respostas foram vagas, informando apenas que a não adoção do procedimento correto leva a acidentes, sem maior detalhamento desses possíveis acidentes.
Um baixo percentual de graduandos de Bacharelado alcançou um resultado satisfatório, sendo a média obtida igual a 24,8.
As principais fontes de enganos que levaram a esse resultado insatisfatório foram : - Item ‘ a ’ : a não explicação da função arrefecedora da água. - Item ‘ b ’ : a não compreensão do volume acumulado de gás, fato que levou a interpretações equivocadas das possíveis consequências. - Item ‘ c ’ : a não compreensão do experimento realizado. - Uma abordagem generalista em todos os itens de que a não adoção do procedimento correto leva a acidentes.
De forma geral, as respostas que receberam as notas mais baixas continham acerto de um dos itens, geralmente o item ‘ a ’ ou o item ‘ b ’.
As respostas medianas apontaram a compreensão de dois dos itens, novamente com destaque para os itens ‘ a ’ e ‘ b ’.
Já as respostas boas foram aquelas nas quais os concluintes conseguiram desenvolver respostas satisfatórias nos três itens, seja construindo um texto correto para os três itens ou ainda correto em dois itens e parcialmente correto em um terceiro.
Como ocorreu nas outras questões discursivas, um reduzido percentual de graduandos conseguiu desenvolver uma resposta satisfatória.
A questão abordava o tema segurança, buscando a discussão e interpretação de situações de laboratório.
Além de possíveis lacunas na formação inicial dos formandos, também foi possível observar a recorrente dificuldade de expressão do conhecimento na construção de um texto discursivo.
Comparando -se os resultados dos concluintes de Bacharelado nas três questões discursivas, observou -se que as médias foram baixas, variando entre 19 e 33.
Esse desempenho muito aquém do esperado pode ser um indicativo da necessidade de uma melhor abordagem dos conteúdos nos respectivos cursos de graduação.
Questão 5 – Licenciatura Essa questão foi específica para os graduandos de Licenciatura.
Os conteúdos abordados envolviam a técnica de ensino-aprendizagem baseada em analogias, uma estratégia comumente empregada pelos professores de Química.
Metodologias de Ensino de Química são fundamentais para o Licenciando e estão alinhadas às Diretrizes que norteiam os cursos de graduação.
O nível de profundidade exigido pode ser considerado médio.
A questão apresentava um enunciado contextualizado, claro e com linguagem adequada, sem possibilidade de interpretações diferentes da esperada.
O padrão de respostas apresentou -se claro e abrangente, possibilitando a pontuação para respostas com desenvolvimento parcial.
Apesar de a questão não apresentar subitens, o enunciado solicitava uma justifica a favor e outra contrária ao uso da metodologia.
Embora não apresentasse subitens, uma vez que o enunciado pedia uma justificativa favorável e uma desfavorável ao uso de analogias, o estudante deveria contemplar duas argumentações distintas.
Em relação às justificativas para o uso de analogias, a maioria das respostas corretas manteve -se dentro do padrão proposto, afirmando que a aproximação do conhecimento do estudante dos conhecimentos específicos e a menor abstração a partir da visualização do modelo são aspectos positivos no ensino.
Quanto às justificativas para o não uso de analogias, a maioria das respostas baseou -se no fato de que as diferenças entre o objeto análogo e o objeto alvo podem levar a erros conceituais.
Dentre as respostas insatisfatórias, as mais comuns apontavam que o modelo de Thomson não reproduz a realidade, explicando modelos posteriores como o de Bohr e o da mecânica quântica.
Apesar de não ser incorreto, este argumento não foi considerado, pois o enunciado buscava a discussão do uso de analogias na abordagem de modelos atômicos e não a comparação entre os diferentes modelos atômicos propostos.
O desempenho dos graduandos de Licenciatura não foi satisfatório, sendo a média obtida igual a 17,1.
Não houve um tipo predominante de erro, mas foi possível constatar que os graduandos de Licenciatura não tinham conhecimento sobre o tema ou que não conseguiram construir uma linha de raciocínio completa em suas respostas.
De forma geral, as respostas que obtiveram as notas mais baixas estiveram relacionadas à incapacidade de construção de um argumento lógico, seja favorável ou desfavorável ao uso de analogias no ensino.
As respostas medianas representaram a compreensão de uma das duas justificativas solicitadas.
Já as respostas boas foram aquelas nas quais os estudantes conseguiram desenvolver respostas satisfatórias para justificar, tanto favoravelmente quanto desfavoravelmente, o uso de analogias no ensino.
Um reduzido percentual de graduandos conseguiu desenvolver uma resposta satisfatória para uma questão envolvendo uma estratégia didática, bem conhecida na área, para o ensino de modelos, no caso específico da questão, modelos atômicos.
Além de possíveis lacunas na formação inicial dos graduandos de Licenciatura, foi possível observar, também aqui, a dificuldade de expressão do conhecimento por meio de um texto discursivo.
Comparando -se os resultados dos graduandos de Licenciatura nas três questões discursivas, observou -se que as médias foram muito baixas, sendo 13,0, 13,1 e 17,1, respectivamente.
O desempenho bem aquém do desejado pode ser um indicativo da necessidade de uma melhor abordagem dos conteúdos relacionados nos respectivos cursos de graduação. 6.3.3 Considerações Finais Analisando -se os resultados globais das questões discursivas, podem -se \textbf{traçar} algumas considerações que venham a contribuir para o desenvolvimento do ensino superior na Área de Química.
A primeira consideração relaciona -se à baixa média alcançada pelos graduandos em todas as respostas.
As médias variaram de 13 e 33,5, o que indica homogeneidade nos resultados, em relação tanto aos formandos do Bacharelado quanto os de Licenciatura, e em relação a todos os conteúdos abordados na prova.
Sem dúvida, as respostas apontam falhas na formação dos graduandos, falhas essas que vão desde à compreensão de conceitos básicos, como equilíbrio químico, cinética química e termoquímica, até à incapacidade de elaboração de um texto argumentativo.
Outra consideração importante relaciona -se ao fato de que, nas duas questões comuns aos graduandos de Bacharelado e de Licenciatura, os bacharéis obtiveram notas ligeiramente superiores, fato que pode ser um indicativo de que a formação inicial em Licenciatura em Química demanda mais cuidados e investimentos.

% matched lemmas: traçar
\end{textsample}
